% ============================================================
%  Cl(0,2) =~ H Two-Observable Decomposition --- Paper 5
%  Todd M. Firestone
%  DRAFT RevA --- not for circulation
%  Typeset for arXiv submission
%  Document class: article + geometry (generic, arXiv-compatible)
%
%  RevA CHANGELOG (2026-06-03):
%    - Removed [P3] (FirestoneP3) cite from §1.1 list
%    - Removed [P3] bibentry from bibliography
%    - Q3 flag closed (Bo entry was complete; flag was stale)
%    - [P4] citations retained in full (27 occurrences)
%
%  FLAGS USED IN THIS FILE:
%    % NOTE:    --- placeholder or missing content
%    % REVIEW:  --- internal inconsistency; author decision required
%    % ASK:     --- ambiguity that cannot be resolved without author input
%    % PENDING: --- content not yet finalized
%
%  STRUCTURAL DECISIONS (per author, this session):
%    Q1: Section-prefixed equation numbering (X.1), (X.2), ...;
%        original-draft gaps closed by re-numbering.
%    Q2: Appendix A starts at A.1 (original draft's A.2 -> A.1, etc.);
%        original A.1 placeholder is absent and was unreferenced.
%    Q3: CLOSED. [Bo] bibliographic entry is present and complete
%        (DOI + arXiv). The earlier NOTE flag was stale.
%    Q4: article + geometry, formatting aligned with Paper 3 draft.
%
%    Theorem-like results numbered by section, independent counters
%    (Theorem 3.1, Lemma 2.1, Proposition 3.1, Corollary 6.1, ...).
%    This re-numbering supersedes the cross-paper sequential numbering
%    in the original draft and resolves the original Cor 4 / Cor 5
%    document-order anomaly.
% ============================================================

\documentclass[12pt]{article}

% --- Encoding & fonts ---
\usepackage[utf8]{inputenc}
\usepackage[T1]{fontenc}
\usepackage{lmodern}
\usepackage{microtype}

% --- Page geometry (arXiv standard margins) ---
\usepackage[top=1in, bottom=1in, left=1.25in, right=1.25in]{geometry}

% --- Mathematics ---
\usepackage{amsmath}
\usepackage{amssymb}
\usepackage{amsthm}
\usepackage{mathtools}

% --- Tables ---
\usepackage{booktabs}
\usepackage{array}
\usepackage{multirow}
\usepackage{longtable}

% --- Graphics & color ---
\usepackage{graphicx}
\usepackage{xcolor}

% --- Citations: numbered, sorted, compressed [1,2,3] style ---
\usepackage[numbers,sort&compress]{natbib}
\bibliographystyle{unsrtnat}

% --- Hyperlinks (arXiv-safe) ---
\usepackage[
colorlinks=true,
linkcolor=blue!60!black,
citecolor=blue!60!black,
urlcolor=blue!60!black,
pdftitle={Two-Observable Decomposition in SU(2) Lattice Gauge
	Theory: the Cl(0,2) = H Framework},
pdfauthor={Todd M. Firestone}
]{hyperref}

% --- Equation numbering: section-prefixed (X.1), (X.2), (X.3)... ---
% (per author decision; section-prefixed style with gaps closed)
\numberwithin{equation}{section}

% ============================================================
%  THEOREM ENVIRONMENTS
%  Section-prefixed numbering; independent counters per result type
%  (so Theorem, Lemma, Proposition, Corollary each have their own
%  X.1, X.2, ... sequence within each section).
% ============================================================
\theoremstyle{plain}
\newtheorem{theorem}{Theorem}[section]
\newtheorem{lemma}{Lemma}[section]
\newtheorem{proposition}{Proposition}[section]
\newtheorem{corollary}{Corollary}[section]
\newtheorem*{corelemma}{Core Lemma}   % unnumbered (sec:core-lemma)

\theoremstyle{definition}
\newtheorem{definition}[theorem]{Definition}

\theoremstyle{remark}
\newtheorem{remark}{Remark}[section]

% ============================================================
%  CUSTOM MACROS
% ============================================================

% --- Inherited from Papers 1, 2, 3 ---
\newcommand{\cl}{\mathrm{Cl}}              % Clifford algebra: \cl(0,2)
\newcommand{\su}{\mathrm{SU}}              % SU group: \su(2)
\newcommand{\uu}{\mathrm{U}}               % U(1) groups
\newcommand{\tr}{\mathrm{Tr}}              % Trace
\newcommand{\half}{\tfrac{1}{2}}
\newcommand{\R}{\mathbb{R}}                % real numbers
\newcommand{\Cplx}{\mathbb{C}}             % complex numbers
\newcommand{\HH}{\mathbb{H}}               % quaternions
\newcommand{\OO}{\mathbb{O}}               % octonions
\newcommand{\Dix}{\R\otimes\Cplx\otimes\HH\otimes\OO}  % Dixon chain

% --- New in Paper 5 ---
\newcommand{\ihat}{\hat{\imath}}           % Cl(0,2) basis vector i-hat
\newcommand{\jhat}{\hat{\jmath}}           % Cl(0,2) basis vector j-hat
\newcommand{\Cjhat}{\Cplx\jhat}            % right C-module of b*jhat
\newcommand{\Gdiag}{G_{\mathrm{diag}}}     % diagonal correlator
\newcommand{\Gcross}{G_{\mathrm{cross}}}   % cross (b-channel) correlator
\newcommand{\Gfull}{G_{\mathrm{full}}}     % full correlator
\newcommand{\Dcross}{\Delta_{\mathrm{cross}}}  % residual cross term
% The superscript ``LW'' on P^{LW}_\pm denotes ``Left-Weak'',
% following the Paper 4 convention for the T_3-eigenspace projectors
% of the SU(2)_L electroweak sector. The macro is retained as \PLW
% for continuity with Paper 4.
\newcommand{\PLW}{P^{\mathrm{LW}}}         % T_3-eigenspace projector (Left-Weak)
\newcommand{\Pip}{\Pi_{+}}                 % rank-1 T_3 eigenprojector (+)
\newcommand{\Pim}{\Pi_{-}}                 % rank-1 T_3 eigenprojector (-)
\newcommand{\Pir}{\Pi_{r}}                 % rank-1 T_3 eigenprojector (r)

% Real-part operator: amsmath/amssymb's \Re is the Fraktur symbol;
% we deliberately do not redefine it. Use \operatorname{Re} below for
% the real part of a bilinear / matrix element, consistent with
% physics-paper convention.

% ============================================================
%  DRAFT HEADER
% ============================================================
% Running header: "Cl(0,2) Two-Observable" -- retained as-is.
% (Paper 3 is no longer active; the parallel header style is moot.)
\usepackage{fancyhdr}
\pagestyle{fancy}
\fancyhf{}
\rhead{\textit{$\cl(0,2)$ Two-Observable
		$\cdot$ DRAFT --- not for circulation}}
\rfoot{\thepage}
\renewcommand{\headrulewidth}{0.4pt}

% ============================================================
\begin{document}
	% ============================================================
	
	% --- Title block ---
	\title{%
		\textbf{Two-Observable Decomposition in
			\texorpdfstring{$\su(2)$}{SU(2)} Lattice Gauge Theory:}\\[4pt]
		\textbf{the \texorpdfstring{$\cl(0,2)\cong\HH$}{Cl(0,2) = H}
			Framework}%
	}
	
	\author{Todd M.\ Firestone\\[2pt]
		\textit{\small }}
	
	% NOTE: Journal target "Advances in Applied Clifford Algebras" and the
	% series identifier "Paper 5 of the R(x)C(x)H(x)O Series" appear in the
	% original draft beneath the byline. They are intentionally not placed
	% in the LaTeX title block here; author to confirm whether they should
	% appear as \thanks{} footnotes, as \date{} content, or remain in
	% comments only until journal submission.
	\date{}
	
	\maketitle
	\thispagestyle{fancy}
	
	% ============================================================
	\begin{abstract}
		% ============================================================
		We present a two-observable decomposition of the $\su(2)$ temporal
		gauge-link correlator within the $\cl(0,2)\cong\HH$ algebraic
		framework.  The $\su(2)$ parametrisation $U = a\cdot I_2 + b\cdot J$,
		with $a = (U)_{00}\in\Cplx$ and $b = (U)_{01}\in\Cjhat$ (a right
		$\Cplx$-module), induces a natural channel decomposition of the
		full correlator $\Gfull(t)$ into two bilinears:
		$\Gdiag(t) = 8\sum \operatorname{Re}[a_t\bar{a}_0]$, probing the
		diagonal ($T_3$ eigenspace) sector, and
		$\Gcross(t) = 8\sum \operatorname{Re}[b_t\bar{b}_0]$, probing the
		off-diagonal sector.  This identification resolves a structural
		gap of Paper~4 of this series: the cross term $\Dcross$, previously
		defined subtractionally, is exactly the $b$-channel bilinear.
		
		Theorem~\ref{thm:two-obs-completeness}
		(Sec.~\ref{sec:theorem-completeness}) establishes
		$\Gdiag(t) + \Gcross(t) = \Gfull(t)$ exactly for every $\su(2)$
		field configuration and every timeslice~$t$.
		Corollary~\ref{cor:two-obs-decomp} (Two-Observable Decomposition,
		Sec.~\ref{sec:cor-two-obs-decomp}) assembles this with the channel
		faithfulness results --- $a$-channel (Proposition~1 of Paper~4)
		and $b$-channel (Proposition~\ref{prop:off-diag-collapse},
		Sec.~\ref{sec:off-diag-collapse}) --- into a complete statement of
		the exact algebraic decomposition.
		Proposition~\ref{prop:rank1-univ} (Sec.~\ref{sec:rank1-univ})
		proves that the rank-1 collapse formula
		$\Pir\,U\,\Pir = (U)_{rr}\,\Pir$, which underlies the channel
		separation, is intrinsic to the fundamental $\su(2)$ representation
		and holds for every tensor-product embedding of $U$ into $\Cplx^8$.
		
		In the weak-coupling expansion, $\Gdiag$ probes the neutral gauge
		boson field $A^3$ and $\Gcross$ probes the charged gauge boson
		fields $W^{\pm}$.
		Proposition~\ref{prop:b-conduit} (Sec.~\ref{sec:b-conduit})
		establishes $\Gcross$ as a necessary algebraic conduit for all
		charged-current coupling.
		Proposition~\ref{prop:gauge-noninv} (Sec.~\ref{sec:gauge-noninv})
		proves $\Gcross$ is not gauge-invariant; under Landau gauge fixing,
		\[
		\Dcross(t)\big|_{\mathrm{Landau}}\big/\,8V
		\;\approx\; +(g^2/2)\,C_{W^{\pm}}(t)
		\quad\text{at } O(g^2),
		\]
		providing a formal connection to the zero-spatial-momentum
		$W^{\pm}$ temporal correlator.
		
		All results are verified by 51 numerical checks (D1--D18), each
		passing at tolerance $\leq 10^{-13}$ or better.  The framework
		does not distinguish $\su(2)_L$ from $\su(2)_R$, which requires
		external fermionic input, and makes no claims about chirality
		selection, $W/Z$ mass ratios, or the Higgs mechanism.
	\end{abstract}
	
	% ============================================================
	\section{Introduction}
	\label{sec:intro}
	% ============================================================
	
	\subsection{Context}
	\label{sec:intro-context}
	
	The present work continues a programme of embedding lattice $\su(2)$
	gauge theory within the division algebra hierarchy $\Dix$, exploiting
	the Clifford algebra isomorphism $\cl(0,2)\cong\HH$ as the natural
	algebraic setting for the electroweak $\su(2)$ sector
	\cite{FirestoneP1,FirestoneP2,FirestoneP4}.
	Paper~4~\cite{FirestoneP4} introduced the $T_3$-eigenspace projectors
	$\PLW_\pm$ and established two central results for the $\su(2)$
	temporal gauge-link correlator.  Proposition~1 of
	Paper~4~\cite{FirestoneP4}, \S4.2, proved that the diagonal bilinear
	collapses to a rank-1 form:
	% ASK: This rank-1 collapse formula (and the off-diagonal collapse
	% formula previewed in §1.2 below) are restatements of upstream results
	% and are unnumbered in the original draft. Promoted display equations
	% may optionally be numbered; left unnumbered here to match the draft.
	% Author to confirm.
	\[
	\operatorname{Re}\,\tr\!\left[\PLW_r\, E_t\,\PLW_r\, E_0^{\dagger}\right]
	\;=\; 4\operatorname{Re}[a_t \bar{a}_0],
	\]
	where $a_t = (U_t)_{00}\in\Cplx$ is the top-left matrix element of
	the $\su(2)$ link variable $U_t$.  Corollary~1 of
	Paper~4~\cite{FirestoneP4} then established that $G_+(t) = G_-(t)$
	identically for every $U\in\su(2)$, so the two eigenspace correlators
	carry identical information and combine to give the diagonal
	correlator $\Gdiag(t)$.
	
	Paper~4 further decomposed the full correlator as
	$\Gfull(t) = \Gdiag(t) + \Dcross(t)$, with the cross term $\Dcross$
	defined in Paper~4~\cite{FirestoneP4}, \S4.3, as a residual:
	$\Dcross(t) = \Gfull(t) - \Gdiag(t)$.  This compact identity was
	verified numerically across four thermalized ensembles
	($\beta_{\mathrm{eff}}\in\{5.0,\,6.0,\,9.0,\,12.0\}$, $N=6$, 20
	configurations) to within
	$|\Gdiag + \Dcross - \Gfull|\leq 8.9\times 10^{-16}$
	\cite{FirestoneP4}.  The ratio $\Dcross/\Gfull$ was observed to lie
	in the range $[0.72,\,2.33]$ across ensembles---a striking $O(1)$
	contribution---and $\Gcross$ was found to be negative in 37--62\% of
	(configuration,\,$t$) pairs~\cite{FirestoneP4}.
	
	Despite this quantitative success, a structural gap remained.  The
	definition of $\Dcross$ in Paper~4 was purely subtractional: it named
	a quantity without providing an independent algebraic generator.
	The question left open was: what is $\Gcross$, exactly, in terms of
	the fundamental $\su(2)$ matrix elements?  Answering this question
	is the task of the present paper.

	\medskip\noindent
	\textbf{What is standard and what is new.}
	The parametrisation $U = a\cdot I_2 + b\cdot J$ with $|a|^2+|b|^2=1$
	is the standard quaternionic form of an $\su(2)$ matrix; the split of
	$\tr[E_t E_0^{\dagger}]$ into $a$- and $b$-channel bilinears follows
	immediately from expanding the matrix product and is not claimed as
	new here.  What is new in this paper is threefold:
	(i)~the \emph{projector formulation} --- identifying the eigenspace
	projectors $\PLW_\pm$ as the algebraic objects that select the $a$-
	and $b$-channels within the $\cl(0,2)\cong\HH$ framework;
	(ii)~the \emph{Paper~4 gap closure} --- proving that the
	subtractionally defined $\Dcross$ equals exactly the $b$-channel
	bilinear $8\sum\operatorname{Re}[b_t\bar{b}_0]$, supplying it with
	an independent algebraic generator for the first time;
	and (iii)~the \emph{Landau-gauge propagator identification} --- the
	formal $O(g^2)$ connection between $\Gcross$ and the
	zero-spatial-momentum $W^{\pm}$ temporal correlator.
	
	\subsection{This Paper}
	\label{sec:intro-this-paper}
	
	The answer is provided by the off-diagonal rank-1 collapse.  Writing
	the $\su(2)$ link variable in the standard parametrisation
	$U = a\cdot I_2 + b\cdot J$, where $a = (U)_{00}\in\Cplx$
	(the $a$-channel) and $b = (U)_{01}\in\Cplx$ (the $b$-channel),
	with $|a|^2 + |b|^2 = 1$, the central new result of this paper
	--- the projector identification of the $b$-channel ---
	previews Proposition~\ref{prop:off-diag-collapse}:
	
	\medskip
	\noindent\emph{Proposition~\ref{prop:off-diag-collapse}
		(Off-Diagonal Rank-1 Collapse).}\enspace
	For every $U_t,\,U_0\in\su(2)$ with $E = I_2\otimes U\otimes I_2$,
	\[
	\operatorname{Re}\,\tr\!\left[\PLW_+\, E_t\,\PLW_-\, E_0^{\dagger}\right]
	\;=\; 4\operatorname{Re}[b_t \bar{b}_0],
	\]
	where $b_t = (U_t)_{01}$.  \emph{The new content is the projector
	identification: $\PLW_\pm$ selects the $b$-channel, giving $\Dcross$
	an independent algebraic generator for the first time.}
	
	\medskip
	From Proposition~\ref{prop:off-diag-collapse}, together with
	Corollary~\ref{cor:orientation-sym}
	(Sec.~\ref{sec:orientation-sym}) and Proposition~1 of
	Paper~4~\cite{FirestoneP4}, the following results are established:
	
	\begin{itemize}
		\item
		\textbf{Theorem~\ref{thm:two-obs-completeness}}
		(Sec.~\ref{sec:theorem-completeness}):
		$\Gdiag(t) + \Gcross(t) = \Gfull(t)$ holds \emph{exactly} for
		every $\su(2)$ gauge configuration and every timeslice $t$,
		where $\Gcross(t) = 8\sum\operatorname{Re}[b_t \bar{b}_0]$.
		Theorem~\ref{thm:two-obs-completeness} simultaneously confirms
		that $\Gcross \equiv \Dcross$.
		
		\item
		\textbf{Corollary~\ref{cor:two-obs-decomp}}
		(Sec.~\ref{sec:cor-two-obs-decomp}, Two-Observable Decomposition):
		the pair $\{\Gdiag,\Gcross\}$ constitutes an exact algebraic
		decomposition of $\Gfull$ in the $\cl(0,2)\cong\HH$ framework.
		
		\item
		\textbf{Proposition~\ref{prop:rank1-univ}}
		(Sec.~\ref{sec:rank1-univ}): the rank-1 collapse is intrinsic to
		the fundamental $\su(2)$ representation, independent of the
		specific tensor-product embedding.
		
		\item
		\textbf{Proposition~\ref{prop:b-conduit}}
		(Sec.~\ref{sec:b-conduit}): the $b$-channel is a necessary
		algebraic conduit for any gauge-link correlator sensitive to
		$W^{\pm}$ at $O(g)$.
		
		\item
		\textbf{Proposition~\ref{prop:gauge-noninv}}
		(Sec.~\ref{sec:gauge-noninv}): $\Dcross$ is not gauge-invariant.
		In Landau gauge,
		$\Dcross(t)/8V \approx +(g^2/2)\,C_{W^{\pm}}(t)$ at $O(g^2)$.
	\end{itemize}
	
	The bridge between Paper~4 and Paper~5 is explicit:
	$\Gcross(t) \equiv \Dcross(t)$, verified numerically to
	$\leq 1.14\times 10^{-13}$ across all four ensembles
	(checks D7a--D7c, D8\,$\times\,4\beta$).
	
	\subsection{Scope Limits}
	\label{sec:intro-scope}
	
	This paper works entirely within the pure $\su(2)$ gauge sector.
	The following topics lie strictly outside its scope and are not
	addressed:
	
	\begin{enumerate}
		\item
		\emph{No Higgs mechanism or electroweak symmetry breaking.}
		The gauge fields remain massless throughout; no scalar sector
		is introduced.
		
		\item
		\emph{No chiral fermions.}  The Nielsen--Ninomiya
		obstruction~\cite{NielsenNinomiya1981a,NielsenNinomiya1981b}
		is not overcome.  Paper~5 contains no fermion action and makes
		no claim of implementing chiral symmetry on the lattice.
		
		\item
		\emph{No $\su(2)_L$/$\su(2)_R$ distinction.}  The algebraic
		framework identifies $\Gcross$ as a $b$-channel probe, but does
		not select between left- and right-handed couplings.
		\emph{The embedding does not, by itself, distinguish
			$\su(2)_L$ from $\su(2)_R$; this requires external fermionic
			input.}
		
		\item
		\emph{No dynamical symmetry-breaking signals.}  Confinement and
		chiral condensate are not discussed; the ensembles are
		pure-gauge and thermalized in the absence of matter fields.
	\end{enumerate}
	
	\subsection{Paper Organisation}
	\label{sec:intro-organization}
	
	Section~\ref{sec:framework} establishes the algebraic framework:
	the $\cl(0,2)\cong\HH$ decomposition, the right $\Cplx$-module
	structure of the $b$-channel (Lemma~\ref{lem:right-Cplx-module}),
	and the $\su(2)$ parametrisation.
	Section~\ref{sec:two-obs-decomp-section} proves
	Proposition~\ref{prop:off-diag-collapse},
	Corollary~\ref{cor:orientation-sym}, and
	Theorem~\ref{thm:two-obs-completeness}, establishing the
	two-observable decomposition algebraically and numerically.
	Section~\ref{sec:physical-interp} interprets the $a$- and
	$b$-channels physically through the weak-coupling expansion.
	Section~\ref{sec:embedding-univ} proves embedding universality
	(Proposition~\ref{prop:rank1-univ}) and the universal $b$-channel
	blindness of diagonal projectors (Lemma~\ref{lem:b-blindness}).
	Section~\ref{sec:two-obs-structure} assembles the two-observable
	decomposition into Corollary~\ref{cor:two-obs-decomp}.
	Section~\ref{sec:SU2-LR} examines the connection to $\su(2)_L$
	through channel assignment (Proposition~\ref{prop:b-conduit}) and
	the conditional prediction for fermionic observables.
	Section~\ref{sec:gauge-properties} proves gauge non-invariance of
	$\Dcross$ (Proposition~\ref{prop:gauge-noninv},
	Lemma~\ref{lem:b-transform}) and develops the Landau gauge
	prescription.
	Section~\ref{sec:discussion} synthesises the results in discussion,
	including five conclusions and seven open questions.
	Section~\ref{sec:conclusions} presents the conclusions.
	
	% ============================================================
	\section{Algebraic Framework: $a$/$b$-Channel Decomposition}
	\label{sec:framework}
	% ============================================================
	
	\subsection{The \texorpdfstring{$\cl(0,2)\cong\HH$}{Cl(0,2) = H}
		Decomposition}
	\label{sec:Cl02-decomp}
	
	The Clifford algebra $\cl(0,2)$ is generated by two anticommuting
	orthonormal vectors $\ihat,\,\jhat$ satisfying
	$\ihat^{\,2} = -1$, $\jhat^{\,2} = -1$, and
	$\ihat\,\jhat = -\jhat\,\ihat$.  The product $\ihat\jhat$ squares to
	$-1$ and anticommutes with both $\ihat$ and $\jhat$, giving the
	dimension count $1 + 2 + 1 = 4$ graded basis elements
	\[
	\{1,\;\ihat,\;\jhat,\;\ihat\jhat\}.
	\]
	As an associative algebra over $\R$, $\cl(0,2)$ is isomorphic to the
	quaternion algebra $\HH$ via the identification
	\[
	1 \;\leftrightarrow\; 1, \qquad
	\ihat \;\leftrightarrow\; i, \qquad
	\jhat \;\leftrightarrow\; j, \qquad
	\ihat\jhat \;\leftrightarrow\; k.
	\]
	The complex subalgebra $\Cplx\subset\cl(0,2)$ is the algebra of
	elements of the form $\alpha + \beta\,\ihat$ with $\alpha,\beta\in\R$;
	the imaginary unit of this $\Cplx$ is identified with $\ihat$.
	
	The decomposition relevant for the present paper is the right-$\Cplx$
	module direct sum
	\begin{equation}
		\cl(0,2) \;=\; \Cplx \,\oplus\, \Cjhat,
		\label{eq:Cl02-decomp}
	\end{equation}
	where $\Cjhat$ is the right $\Cplx$-module of elements $b\,\jhat$ with
	$b\in\Cplx$.  The action of $\Cplx$ on this module is from the right:
	$(b\,\jhat)\cdot c = b(\jhat c) = b\,\bar c\,\jhat$ for any
	$b,c\in\Cplx$, using the anticommutation $\jhat\,c = \bar c\,\jhat$.
	The decomposition~\eqref{eq:Cl02-decomp} is the algebraic foundation
	of the two-observable decomposition.
	
	% ------------------------------------------------------------
	\subsection{\texorpdfstring{$\Cjhat$}{Cj-hat} is a Right
		\texorpdfstring{$\Cplx$}{C}-Module}
	\label{sec:right-Cplx-module}
	% ------------------------------------------------------------
	
	The structure $\Cjhat$ is precisely a right $\Cplx$-module, not a
	left $\HH$-module.
	
	\begin{lemma}\label{lem:right-Cplx-module}
		For any $b\in\Cplx$ and $c\in\Cplx$, defining the right action of
		$\Cplx$ on the formal expression $b\,\jhat$ by
		\[
		(b\,\jhat)\cdot c \;:=\; b(\jhat c) \;=\; b\,\bar c\,\jhat
		\qquad (\text{using } \jhat\,c = \bar c\,\jhat),
		\]
		the set $\Cjhat := \{b\,\jhat : b \in \Cplx\}$ forms a one-dimensional
		right $\Cplx$-module under this action.
	\end{lemma}
	
	Properties RM1--RM5 are verified by direct expansion using the
	anticommutation rule $\jhat\,c = \bar c\,\jhat$ (checks D2a--d).
	The verification is summarised below:
	
	\begin{center}
		\small
		\setlength{\tabcolsep}{6pt}
		\begin{tabular}{@{}lll@{}}
			\toprule
			\textbf{Property} & \textbf{Statement} & \textbf{Verified}\\
			\midrule
			RM1 & $\bigl((b_1 + b_2)\jhat\bigr)\cdot c
			= (b_1\jhat\cdot c) + (b_2\jhat\cdot c)$
			& \checkmark \\[2pt]
			RM2 & $(b\,\jhat)\cdot(c_1 + c_2)
			= (b\,\jhat\cdot c_1) + (b\,\jhat\cdot c_2)$
			& \checkmark \\[2pt]
			RM3 & $\bigl((b\,\jhat)\cdot c_1\bigr)\cdot c_2
			= (b\,\jhat)\cdot(c_1 c_2)$
			& \checkmark \\[2pt]
			RM4 & $(b\,\jhat)\cdot 1 = b\,\jhat$
			& \checkmark \\[2pt]
			RM5 & One-dimensional: spanned by $\{\jhat\}$ over $\Cplx$
			& \checkmark \\
			\bottomrule
		\end{tabular}
	\end{center}
	
	\noindent
	The verification establishes that $\Cjhat$ is the formal
	one-dimensional right $\Cplx$-module, isomorphic to $\Cplx$ as a
	$\Cplx$-vector space but with the scalar action carrying a complex
	conjugation.
	
	\medskip\noindent
	\textbf{This module structure is preserved exactly under the
		parametrisation $U = a\cdot I_2 + b\cdot J$ in
		Sec.~\ref{sec:SU2-param}.}  The $b$-channel of the $\su(2)$
	decomposition inherits the right $\Cplx$-module structure: $b\in\Cplx$
	is the module coefficient, $J$ is the analog of $\jhat$.
	
	\medskip\noindent
	\emph{Per the standing correction:} this is a right $\Cplx$-module,
	\emph{not} a left $\HH$-module.  Earlier references to a left
	$\HH$-module structure are incorrect; the correct algebraic identity
	is the right $\Cplx$-module decomposition~\eqref{eq:Cl02-decomp}.
	
	% ------------------------------------------------------------
	\subsection{\texorpdfstring{$\su(2)$}{SU(2)} Parametrisation}
	\label{sec:SU2-param}
	% ------------------------------------------------------------
	
	Every $\su(2)$ matrix admits the parametrisation
	\begin{equation}
		U \;=\;
		\begin{pmatrix} a & b \\ -b^* & a^* \end{pmatrix}
		\;=\; a\cdot I_2 + b\cdot J
		\qquad (|a|^2 + |b|^2 = 1),
		\label{eq:SU2-param}
	\end{equation}
	where:
	\begin{itemize}
		\item $a = (U)_{00} \in \Cplx$, called the \emph{$a$-channel};
		\item $b = (U)_{01} \in \Cplx$, called the \emph{$b$-channel}.
	\end{itemize}
	The matrix $J$ in~\eqref{eq:SU2-param} is the $2\times 2$
	representation of the quaternion unit $\jhat$:
	\begin{equation}
		J \;=\; i\sigma_2 \;=\;
		\begin{pmatrix} 0 & 1 \\ -1 & 0 \end{pmatrix},
		\label{eq:J-def}
	\end{equation}
	satisfying $J^2 = -I_2$ and $J^{\dagger} = -J$, consistent with
	the Clifford relation $\jhat^2 = -1$.
	
	When the $\cl(0,2)\cong\HH$ decomposition~\eqref{eq:Cl02-decomp} is
	enforced, $b$ is more naturally written as $b\,\jhat\in\Cjhat$, and
	the correlator bilinears acquire the right $\Cplx$-module
	interpretation:
	\[
	\operatorname{Re}[b_t\,\bar{b}_0]
	\;\;\longleftrightarrow\;\;
	\text{``right $\Cplx$-module bilinear in $\Cjhat$''.}
	\]
	This identification is what gives $\Gcross$ its algebraic identity
	as the $b$-channel observable.
	
	The fundamental $\su(2)$ representation in this paper is
	$8$-dimensional, embedded as $U\in M_2(\Cplx) \to
	I_2 \otimes U \otimes I_2 \in M_8(\Cplx)$.  All assertions in this
	paper are independent of the choice of tensor-product ordering; see
	Proposition~\ref{prop:rank1-univ} in Sec.~\ref{sec:rank1-univ} for
	the precise universality statement.
	
	For verification, all bench tests use:
	\begin{itemize}
		\item $N = 6^{3}\times 6$ lattice ($V = 6^{4} = 1296$);
		\item $\beta_{\mathrm{eff}}\in\{5.0,\,6.0,\,9.0,\,12.0\}$;
		\item 20 thermalized configurations per ensemble;
		\item Standard heat-bath update; eigenspace projectors $\PLW_\pm$
		from Paper~4~\cite{FirestoneP4}.
	\end{itemize}
	
	A complete map of the algebraic identifications is given in
	Table~\ref{tab:algebraic-ID}.
	
	\begin{table}[ht]
		\centering
		\caption{Algebraic identifications and physical role assignments
			for the $\cl(0,2)\cong\HH\,\leftrightarrow\,\su(2)\,
			\leftrightarrow\,M_2(\Cplx)$ chain.  The fundamental $\su(2)$
			representation (8-dim) provides the chain; each level inherits
			the right-$\Cplx$-module structure of the $b$-channel.}
		\label{tab:algebraic-ID}
		\footnotesize
		\setlength{\tabcolsep}{5pt}
		\begin{tabular}{@{}llll@{}}
			\toprule
			\textbf{Object} & \textbf{Algebraic ID}
			& \textbf{Matrix Form} & \textbf{Role}\\
			\midrule
			$\ihat$ & imaginary unit & $i\sigma_3$
			& Cartan element ($T_3$ orientation)\\
			$\jhat$ & second generator & $i\sigma_2$
			& off-diagonal swap\\
			$\ihat\jhat$ & pseudoscalar & $-i\sigma_1$
			& volume form\\
			$1$ & identity & $I_2$
			& scalar\\
			$a\in\Cplx$ & $a$-channel & $(U)_{00}$
			& diagonal pair $\{a,a^*\}$\\
			$b\in\Cplx$ & $b$-channel & $(U)_{01}$
			& off-diagonal pair $\{b,-b^*\}$\\
			$\PLW_\pm$ & rank-1 $T_3$-eigenprojector
			& $(I_2 \pm \sigma_3)/2$ & eigenspace projector\\
			$\Pip$ & rank-1 fundamental & $(I_2 + \sigma_3)/2$
			& $a$-channel selector\\
			$\Pim$ & rank-1 fundamental & $(I_2 - \sigma_3)/2$
			& $b$-channel selector\\
			$\Gdiag(t)$ & diagonal observable
			& $\operatorname{Re}\bigl(\tr[\PLW_+ E_t \PLW_+ E_0^{\dagger}]
			+ \tr[\PLW_- E_t \PLW_- E_0^{\dagger}]\bigr)$
			& $a$-channel bilinear\\
			$\Gcross(t)$ & off-diagonal observable
			& $\operatorname{Re}\bigl(\tr[\PLW_+ E_t \PLW_- E_0^{\dagger}]
			+ \tr[\PLW_- E_t \PLW_+ E_0^{\dagger}]\bigr)
			= 8\!\sum\operatorname{Re}[b_t\,\bar{b}_0]$
			\quad(Eq.~\ref{eq:Gcross-def})
			& $b$-channel bilinear\\
			\bottomrule
		\end{tabular}
	\end{table}
	
	% ============================================================
	\section{Off-Diagonal Rank-1 Collapse and Two-Observable Completeness}
	\label{sec:two-obs-decomp-section}
	% ============================================================
	
	This section proves Proposition~\ref{prop:off-diag-collapse} and
	Corollary~\ref{cor:orientation-sym}, then assembles the two-observable
	decomposition Theorem~\ref{thm:two-obs-completeness}.
	
	% ------------------------------------------------------------
	\subsection{Off-Diagonal Rank-1 Collapse}
	\label{sec:off-diag-collapse}
	% ------------------------------------------------------------
	
	\begin{proposition}[Off-Diagonal Rank-1 Collapse]
		\label{prop:off-diag-collapse}
		For every $U_t,\,U_0\in\su(2)$ with $E = I_2\otimes U\otimes I_2$:
		\begin{equation}
			\operatorname{Re}\,\tr\!\left[\PLW_+\, E_t\,\PLW_-\,
			E_0^{\dagger}\right]
			\;=\; 4\operatorname{Re}[b_t\,\bar{b}_0],
			\label{eq:off-diag-collapse}
		\end{equation}
		where $b_t = (U_t)_{01}$.
	\end{proposition}
	
	\begin{proof}
		Recall the rank-1 collapse identity: for any $A\in M_2(\Cplx)$,
		\[
		\Pir\, A\,\Pi_{r'}
		\;=\; (A)_{rr'}\,\Pir\,\Pi_{r'}
		\;=\; (A)_{rr'}\,\delta_{rr'}\,\Pir
		\qquad \text{(no sum).}
		\]
		Applied to $A = E_t = I_2\otimes U_t\otimes I_2$ with
		$\Pip,\,\Pim$ the rank-1 fundamental projectors acting on the
		central $\su(2)$ factor:
		\[
		\PLW_+\, E_t\,\PLW_-
		\;=\; (I_2\otimes U_t\otimes I_2)_{01}\,
		(I_2\otimes\Pip\Pim\otimes I_2)
		\;=\; (U_t)_{01}\cdot\text{(rank-1 vanishing term)}
		\;=\; b_t\;(\Pip\Pim\text{ in the trace structure}).
		\]
		Tracing against $E_0^{\dagger}$:
		\[
		\tr\!\left[\PLW_+\, E_t\,\PLW_-\, E_0^{\dagger}\right]
		\;=\; b_t\cdot(E_0^{\dagger})_{10}
		\;=\; b_t\,\bar{b}_0,
		\]
		using $(E_0^{\dagger})_{10} = ((U_0)_{01})^{*} = \bar{b}_0$.
		Taking the real part:
		\[
		\operatorname{Re}\,\tr\!\left[\PLW_+\, E_t\,\PLW_-\,
		E_0^{\dagger}\right]
		\;=\; \operatorname{Re}[b_t\,\bar{b}_0].
		\]
		The factor of $4$ arises from the $I_2\otimes(\cdot)\otimes I_2$
		tensor product:
		\[
		\operatorname{Re}\,\tr_{8\times 8}\!
		\left[\PLW_+\, E_t\,\PLW_-\, E_0^{\dagger}\right]
		\;=\; 4\operatorname{Re}\,\tr_{2\times 2}\!
		\left[\Pip\, U_t\,\Pim\, U_0^{\dagger}\right]
		\;=\; 4\operatorname{Re}[b_t\,\bar{b}_0]. \qedhere
		\]
	\end{proof}
	
	The verification was performed across all four ensembles using checks
	D1, D2a--d, D3a--d, D4a--h, D5, D6, D7a--c with results listed in
	Table~\ref{tab:prop1-verification}.
	
	
	\begin{table}[ht]
		\centering
		\caption{Verification checks for
			Proposition~\ref{prop:off-diag-collapse}.
			All 18 check IDs (51 instances across D1--D18) pass
			at tolerance $\leq 10^{-13}$ or better.}
		\label{tab:prop1-verification}
		\small
		\setlength{\tabcolsep}{6pt}
		\begin{tabular}{@{}lp{8.7cm}l@{}}
			\toprule
			\textbf{Check} & \textbf{Result} & \textbf{Tolerance}\\
			\midrule
			D1 & Rank-1 collapse identity
			$\Pir\, A\,\Pir = (A)_{rr}\,\Pir$ --- algebraic
			& exact \\[2pt]
			D2a--d & $\jhat\,c = \bar{c}\,\jhat$ for
			$c = 1,\,\ihat,\,\beta,\,\ihat\cdot\beta$
			& $\leq 10^{-15}$ \\[2pt]
			D3a--d & $(b\,\jhat)\cdot c = b\cdot\bar{c}\cdot\jhat$
			verification
			& $\leq 10^{-15}$ \\[2pt]
			D4a--h & Direct rank-1 identity in 8-dim:
			$\Pir\, E\,\Pi_{r'} = \delta_{rr'}\,(M)_{rr}$
			& $\leq 10^{-13}$ \\[2pt]
			D5 & Off-diagonal collapse
			$\Pip\, E\,\Pim = (E)_{01}\,\Pip\Pim$
			& $\leq 10^{-13}$ \\[2pt]
			D6 & Off-diagonal rank-1 collapse identity ($8\times 8$):
			$\tr[\Pip\, E\,\Pim\, E'^{\,\dagger}]
			= 4\,(U)_{01}\,(U')_{01}^{*}$
			& $\leq 10^{-13}$ \\[2pt]
			D7a--c & $\operatorname{Re}\tr[\PLW_+ E_t \PLW_- E_0^{\dagger}]
			- 4\operatorname{Re}[b_t\,\bar{b}_0] = 0$ across ensembles
			& $\leq 10^{-13}$ \\
			\bottomrule
		\end{tabular}
	\end{table}
	
	% ------------------------------------------------------------
	\subsection{Orientation Symmetry of the Off-Diagonal Collapse}
	\label{sec:orientation-sym}
	% ------------------------------------------------------------
	
	\begin{corollary}[Orientation Symmetry]
		\label{cor:orientation-sym}
		For every $U_t,\,U_0\in\su(2)$:
		\begin{equation}
			\operatorname{Re}\,\tr\!\left[\PLW_-\, E_t\,\PLW_+\,
			E_0^{\dagger}\right]
			\;=\; \operatorname{Re}\,\tr\!\left[\PLW_+\, E_t\,\PLW_-\,
			E_0^{\dagger}\right]
			\;=\; 4\operatorname{Re}[b_t\,\bar{b}_0].
			\label{eq:orientation-sym}
		\end{equation}
	\end{corollary}
	
	\begin{proof}
		Both reduce to the same expression upon application of
		Proposition~\ref{prop:off-diag-collapse} to the conjugate pair
		$\{\PLW_+,\,\PLW_-\}$.  The orientation reversal exchanges
		$\tr[X^{\dagger}] = \tr[X]^{*}$ for the real part:
		\[
		\operatorname{Re}\,\tr\!\left[\PLW_-\, E_t\,\PLW_+\,
		E_0^{\dagger}\right]
		\;=\; \operatorname{Re}\,\tr\!\left[\bigl(\PLW_-\, E_t\,\PLW_+\,
		E_0^{\dagger}\bigr)^{\dagger}\right]
		\;=\; \operatorname{Re}\,\tr\!\left[E_0\,\PLW_+\, E_t^{\dagger}\,
		\PLW_-\right].
		\]
		Cycling and using $\PLW_r{}^{2} = \PLW_r$:
		\[
		=\; \operatorname{Re}\,\tr\!\left[\PLW_+\, E_t^{\dagger}\,
		\PLW_-\, E_0\right]
		\;=\; (E_t^{\dagger})_{01}\,(E_0)_{10}^{*}
		\;=\; (\bar{b}_t)(\bar{b}_0)^{*}
		\;=\; \bar{b}_t\,b_0
		\;=\; (b_t\,\bar{b}_0)^{*}.
		\]
		So
		$\operatorname{Re}\,\tr[\PLW_-\, E_t\,\PLW_+\, E_0^{\dagger}]
		= \operatorname{Re}[\bar{b}_t\,b_0]
		= \operatorname{Re}[b_t\,\bar{b}_0]$.
	\end{proof}
	
	The orientation symmetry confirms that the off-diagonal bilinear is a
	single channel (not two), justifying the single $\Gcross$ observable.
	
	% ------------------------------------------------------------
	\subsection{Two-Observable Completeness}
	\label{sec:theorem-completeness}
	% ------------------------------------------------------------
	
	\begin{theorem}[Two-Observable Completeness]
		\label{thm:two-obs-completeness}
		For every $\su(2)$ gauge field configuration and every~$t$:
		\begin{equation}
			\Gdiag(t) + \Gcross(t) \;=\; \Gfull(t),
			\label{eq:two-obs-completeness}
		\end{equation}
		where
		\begin{align}
			\Gdiag(t)  &\;=\; 8\sum
			\operatorname{Re}[a_t\,\bar{a}_0],
			\label{eq:Gdiag-def} \\
			\Gcross(t) &\;=\; 8\sum
			\operatorname{Re}[b_t\,\bar{b}_0],
			\label{eq:Gcross-def} \\
			\Gfull(t)  &\;=\; \sum
			\operatorname{Re}\,\tr\!\left[E_t\, E_0^{\dagger}\right].
			\label{eq:Gfull-def}
		\end{align}
		The sums are over space and over both the (configuration,~$t$)
		pairs.
	\end{theorem}
	
	\begin{proof}
		Expand $E_t E_0^{\dagger}$ using the $\su(2)$
		parametrisation~\eqref{eq:SU2-param}:
		\[
		\tr\!\left[E_t\, E_0^{\dagger}\right]
		\;=\; \tr\!\left[(a_t I_2 + b_t J)\,
		(\bar{a}_0 I_2 - \bar{b}_0 J)\right].
		\]
		Inserting completeness $\Pip + \Pim = I_8$ (or $I_2$ in the
		central factor):
		\[
		E_t\, E_0^{\dagger}
		\;=\; (\Pip + \Pim)\, E_t E_0^{\dagger}\,(\Pip + \Pim)
		\;=\; \Pip E_t E_0^{\dagger}\Pip
		+\Pip E_t E_0^{\dagger}\Pim
		+\Pim E_t E_0^{\dagger}\Pip
		+\Pim E_t E_0^{\dagger}\Pim.
		\]
		Tracing and taking real parts: the two diagonal blocks
		($\Pip\cdots\Pip$, $\Pim\cdots\Pim$) each contribute
		$4\operatorname{Re}[a_t\,\bar{a}_0]$ by the diagonal collapse
		identity (Proposition~1 of Paper~4~\cite{FirestoneP4}), summing to
		$8\operatorname{Re}[a_t\,\bar{a}_0]$.  The two off-diagonal blocks
		($\Pip\cdots\Pim$, $\Pim\cdots\Pip$) each contribute
		$4\operatorname{Re}[b_t\,\bar{b}_0]$ by
		Proposition~\ref{prop:off-diag-collapse} and
		Corollary~\ref{cor:orientation-sym}, summing to
		$8\operatorname{Re}[b_t\,\bar{b}_0]$.  Therefore:
		\[
		\operatorname{Re}\,\tr\!\left[E_t\, E_0^{\dagger}\right]
		\;=\; 8\operatorname{Re}[a_t\,\bar{a}_0]
		+ 8\operatorname{Re}[b_t\,\bar{b}_0].
		\]
		Summing over space gives $\Gfull = \Gdiag + \Gcross$.
	\end{proof}
	
	The simultaneous decomposition was verified numerically:
	check D7a--c gives $|\Gdiag + \Gcross - \Gfull|\leq 1.78\times 10^{-15}$
	across all four ensembles, confirming
	Theorem~\ref{thm:two-obs-completeness} to the precision of the
	simulation; check D8\,$\times\,4\beta$ additionally gives
	$|\Gcross - \Dcross|\leq 1.14\times 10^{-13}$, confirming the bridge
	identification $\Gcross\equiv\Dcross$.
	
	Theorem~\ref{thm:two-obs-completeness} simultaneously establishes
	$\Gcross \equiv \Dcross$: the algebraically defined $\Gcross$ equals
	the subtractionally defined $\Dcross$.  This resolves the structural
	gap of Paper~4~\cite{FirestoneP4}, \S4.3: $\Dcross$ is not merely a
	residual, but the $b$-channel bilinear in the right $\Cplx$-module
	decomposition.
	
	% ============================================================
	\section{Physical Interpretation:
		\texorpdfstring{$W^{3}$}{W3} vs.\ \texorpdfstring{$W^{\pm}$}{W+-}}
	\label{sec:physical-interp}
	% ============================================================
	
	% NOTE: Subsections §4(a)-§4(d) in the original draft used letter-style
	% labels. These have been re-rendered as §4.1-§4.4 for consistency
	% with the numeric subsection style used in all other sections.
	
	\subsection{Weak-Coupling Expansion of the Link Variable}
	\label{sec:weak-coupling-exp}
	
	Setting $U_\mu(x) = \exp\!\bigl(i g\,a_{\mathrm{lat}}\,A^c_\mu T^c\bigr)$
	and expanding to leading order in the coupling $g$, with generators
	$T^c = \sigma^c/2$ in the fundamental representation:
	\begin{align}
		a &= (U)_{00} \;=\; 1 + \frac{ig}{2}\,A^{3}_\mu + O(g^2),
		\label{eq:a-expansion}\\[2pt]
		b &= (U)_{01} \;=\; \frac{ig}{2}\,
		\bigl(A^{1}_\mu - i\,A^{2}_\mu\bigr) + O(g^2),
		\label{eq:b-expansion-naive}\\[2pt]
		b &= (U)_{01} \;=\; \frac{ig}{\sqrt{2}}\,W^{+}_\mu + O(g^2),
		\qquad
		W^{+}_\mu \;=\;
		\frac{A^{1}_\mu - i\,A^{2}_\mu}{\sqrt{2}}.
		\label{eq:b-expansion-W}
	\end{align}
	
	\subsection{Channel-to-Boson Identification}
	\label{sec:channel-boson}
	
	Equations~\eqref{eq:a-expansion}--\eqref{eq:b-expansion-W} establish
	a direct correspondence at $O(g)$:
	\[
	\Gdiag \;\;\longleftrightarrow\;\; A^{3}_\mu
	\quad\text{(neutral gauge boson sector)},
	\qquad
	\Gcross \;\;\longleftrightarrow\;\; W^{\pm}_\mu
	\quad\text{(charged gauge boson sector)}.
	\]
	This identification holds to $O(g)$ in the matrix elements and to
	$O(g^2)$ in the bilinears.  The channel separation---$a$ diagonal,
	$b$ off-diagonal---is an exact property of the $\su(2)$ matrix
	structure and holds for all $g$.
	
	\begin{table}[ht]
		\centering
		\caption{Weak-coupling field content of the $\su(2)$ link variable.
			The third column uses $A^{3}$ for the leading-order neutral
			field.}
		\label{tab:weak-coupling-content}
		\small
		\setlength{\tabcolsep}{8pt}
		\begin{tabular}{@{}llll@{}}
			\toprule
			\textbf{Matrix entry} & \textbf{$\su(2)$ param}
			& \textbf{Leading order in $g$} & \textbf{Gauge boson}\\
			\midrule
			$(U)_{00}$ & $a$
			& $1 + (ig/2)\,A^{3}$ & Neutral ($W^{3}$)\\
			$(U)_{01}$ & $b$
			& $(ig/\sqrt{2})\,W^{+}$ & Charged ($W^{+}$)\\
			$(U)_{10}$ & $-b^{*}$
			& $(ig/\sqrt{2})\,W^{-}$ & Charged ($W^{-}$)\\
			$(U)_{11}$ & $a^{*}$
			& $1 - (ig/2)\,A^{3}$ & Neutral ($W^{3}$)\\
			\bottomrule
		\end{tabular}
	\end{table}
	
	\subsection{\texorpdfstring{$O(g^{2})$}{O(g\^{}2)} Calculation
		of the Channel Bilinears}
	\label{sec:Og2-bilinears}
	
	\begin{align}
		\operatorname{Re}[a_t\,\bar{a}_0]
		&\;=\; 1 + \frac{g^{2}}{4}\,A^{3}_t\,A^{3}_0 + O(g^{3}),
		\label{eq:a-bilinear-Og2}\\[2pt]
		\operatorname{Re}[b_t\,\bar{b}_0]
		&\;=\; \frac{g^{2}}{4}\,
		\bigl(A^{1}_t\,A^{1}_0 + A^{2}_t\,A^{2}_0\bigr) + O(g^{3}).
		\label{eq:b-bilinear-Og2}
	\end{align}
	Equation~\eqref{eq:b-bilinear-Og2} shows that
	$\operatorname{Re}[b_t\,\bar{b}_0]$ is $O(g^{2})$: it vanishes as
	$g\to 0$ (free field).  The measured ratio
	$\Dcross/\Gfull \in [0.72,\,2.33]$ at
	$\beta_{\mathrm{eff}}\in\{5.0,\,6.0,\,9.0,\,12.0\}$
	\cite{FirestoneP4} reflects non-perturbative finite-coupling
	contributions where the $O(g^{2})$ approximation is not quantitatively
	valid.
	
	\begin{remark}[D10b --- Slope of $b$-Channel Weak-Coupling Error]
		\label{rem:D10b-slope}
		The log-log slope of the $b$-channel approximation error vs.\
		$\varepsilon$ is $3.000$, not $2$ (check D10b, flagged and closed).
		This arises because $H^{2}\propto I_{2}$, so $(H^{2})_{01} = 0$,
		and the Taylor expansion of $\sin(\varepsilon\,\|H\|)$ has no
		$\varepsilon^{2}$ term in the off-diagonal entry; the leading
		error is $O(\varepsilon^{3})$.  This is the sin-Taylor argument.
		It is distinct from the BCH/$[T_{+},T_{-}]\propto T_{3}$ argument
		(check D17), which explains channel purity, not slope.
		Slope $=3$ confirms approximation quality; channel identity is
		established algebraically by
		Proposition~\ref{prop:off-diag-collapse},
		Lemma~\ref{lem:b-blindness}, and checks D3c--D3d.
		Checks D11a, D11b, D12b, D13a verify the $a$-channel slope,
		monotone $|R(\beta)|$ at $t=5$, and cold-start
		$\langle|b|^{2}\rangle$.
	\end{remark}
	
	\subsection{What the Channel Identification Does Not Imply}
	\label{sec:not-imply}
	
	\begin{description}
		\item[(OS4.1)] \emph{No $W/Z$ mass ratio.}
		No Weinberg angle $\theta_W$, no $\uu(1)_Y$, no Higgs sector.
		
		\item[(OS4.2)] \emph{No L/R distinction.}
		The $\su(2)$ link variable $U$ transforms identically regardless
		of the handedness of any matter field; chirality is encoded in
		the matter sector only.
		
		\item[(OS4.3)] \emph{No Higgs mechanism.}
		The gauge fields in the present $\su(2)$ pure-gauge theory are
		massless.
		
		\item[(OS4.4)] \emph{No electroweak mixing.}
		The hypercharge gauge group $\uu(1)_Y$ is absent; $A^{3}$ should
		not be identified with the physical $Z$ boson or photon prior to
		mixing.
	\end{description}
	
	% ============================================================
	\section{Embedding Universality}
	\label{sec:embedding-univ}
	% ============================================================
	
	\subsection{The Core Rank-1 Collapse Lemma}
	\label{sec:core-lemma}
	
	\begin{corelemma}[Rank-1 Collapse]
		For any $U\in\mathrm{End}(\Cplx^{2})$ and any rank-1 projector
		$\Pir = |r\rangle\langle r|$:
		\begin{equation}
			\Pir\, U\,\Pir \;=\; (U)_{rr}\cdot\Pir.
			\label{eq:core-lemma}
		\end{equation}
	\end{corelemma}
	
	\begin{proof}
		$\Pir\, U\,\Pir
		= |r\rangle\langle r|\,U\,|r\rangle\langle r|
		= (U)_{rr}\,|r\rangle\langle r|
		= (U)_{rr}\,\Pir.$
	\end{proof}
	
	% ------------------------------------------------------------
	\subsection{Rank-1 Collapse Universality}
	\label{sec:rank1-univ}
	% ------------------------------------------------------------
	
	We consider three natural embeddings of $U\in\su(2)$ into
	$\mathrm{End}(\Cplx^{8})$:
	\begin{align*}
		\text{Case A:} && E &\;=\; U\otimes I_{4},\\
		\text{Case B:} && E &\;=\; I_{4}\otimes U,\\
		\text{Case C:} && E &\;=\; I_{2}\otimes U\otimes I_{2}
		\quad\text{(used in Papers~4 and~5)}.
	\end{align*}
	
	\begin{proposition}[Rank-1 Collapse Universality]
		\label{prop:rank1-univ}
		For each embedding Case A, B, C and for $r\in\{+,-\}$:
		\begin{equation}
			\operatorname{Re}\,\tr\!\left[P_{r}\, E_{t}\, P_{r}\,
			E_{0}^{\dagger}\right]
			\;=\; 4\operatorname{Re}\!\left[(U_t)_{rr}\cdot
			\overline{(U_0)_{rr}}\right].
			\label{eq:rank1-univ}
		\end{equation}
	\end{proposition}
	
	\begin{proof}
		In each case, $P_r$ factors as a tensor product with $\Pir$
		acting on the $U$-component and identity operators on the
		spectator components.  Tensor multiplicativity gives
		\[
		\tr\!\left[P_{r}\, E_{t}\, P_{r}\, E_{0}^{\dagger}\right]
		\;=\; \tr(I_{\mathrm{spec}})\cdot
		\tr\!\left[\Pir\, U_{t}\,\Pir\, U_{0}^{\dagger}\right],
		\]
		where $\tr(I_{\mathrm{spec}}) = 4$ in all three cases.  Applying
		the Core Lemma~\eqref{eq:core-lemma} and evaluating the trace
		gives $4\operatorname{Re}\!\bigl[(U_t)_{rr}\,
		\overline{(U_0)_{rr}}\bigr]$.
	\end{proof}
	
	Verified for Cases A and B by checks D14$^{+}$, D14$^{-}$,
	D15$^{+}$, D15$^{-}$ on 50 random $\su(2)$ pairs each.  Residuals are
	$2.22\times 10^{-16}$ (1 ULP) in all cases.  These are type~(b)
	independent computations: the LHS uses an $8\times 8$ Kronecker-product
	trace, the RHS accesses the $2\times 2$ matrix element directly; the
	two code paths share no intermediate computation.
	
	% ------------------------------------------------------------
	\subsection{Intrinsic Nature of the Decomposition}
	\label{sec:cor-intrinsic}
	% ------------------------------------------------------------
	
	\begin{corollary}[Universality of the Two-Observable Structure]
		\label{cor:intrinsic}
		The rank-1 collapse formula $\Pir\, U\,\Pir = (U)_{rr}\,\Pir$
		---and hence the $a$/$b$-channel separation and the
		two-observable decomposition of
		Theorem~\ref{thm:two-obs-completeness}---is intrinsic to the
		fundamental spin-$1/2$ representation of $\su(2)$.  It is not an
		artifact of the tensor-product embedding $I_{2}\otimes U\otimes
		I_{2}$ adopted in Paper~4.
		\emph{Escape condition:} reducible representations with degenerate
		$T_{3}$ eigenspaces avoid rank-1 collapse; no such representations
		appear in this series.
	\end{corollary}
	
	% ------------------------------------------------------------
	\subsection{\texorpdfstring{$b$}{b}-Channel Blindness is Universal}
	\label{sec:b-blindness}
	% ------------------------------------------------------------
	
	\begin{lemma}[$b$-Channel Blindness Universality]
		\label{lem:b-blindness}
		For every embedding Case A, B, C, the diagonal projector
		correlator
		\[
		\operatorname{Re}\,\tr\!\left[P_{r}\, E_{t}\, P_{r}\,
		E_{0}^{\dagger}\right]
		\;=\; 4\operatorname{Re}\!\left[(U_{t})_{rr}\cdot
		\overline{(U_{0})_{rr}}\right]
		\]
		is independent of $b = (U)_{01}$ for every
		$U_t, U_0\in\su(2)$.
	\end{lemma}
	
	\begin{proof}
		The RHS of~\eqref{eq:rank1-univ} contains only $(U)_{rr} = a$
		(for $r = +$) or $a^{*}$ (for $r = -$), and no occurrence of
		$b = (U)_{01}$.  Independence follows from
		Proposition~\ref{prop:rank1-univ}.
	\end{proof}
	
	Check D16 provides a consistency check over $800$ evaluations,
	finding std $= 0.00\times 10^{+00}$ for the dependence on $b$.  This
	result is guaranteed by construction ($b$ is algebraically absent
	from the formula) and should be read as a consistency check rather
	than an independent verification.  The substantive verification is
	established by Proposition~\ref{prop:rank1-univ} and the Core
	Lemma~\eqref{eq:core-lemma}, and confirmed by the independent
	code-path comparison of D14/D15.
	
	\begin{table}[ht]
		\centering
		\caption{Three-embedding comparison.  Each embedding yields the
			same diagonal collapse formula and the same $b$-blindness,
			differing only in how the spectator dimensions assemble to
			$\tr(I_{\mathrm{spec}}) = 4$.}
		\label{tab:three-embedding}
		\small
		\setlength{\tabcolsep}{6pt}
		\begin{tabular}{@{}cllcc@{}}
			\toprule
			\textbf{Case} & \textbf{Embedding} & \textbf{Formula}
			& \textbf{$b$-blind?} & \textbf{Spectator dim}\\
			\midrule
			A & $U\otimes I_{4}$
			& $4\operatorname{Re}\!\bigl[(U_t)_{rr}\,
			\overline{(U_0)_{rr}}\bigr]$
			& Yes & $\tr(I_{4}) = 4$\\
			B & $I_{4}\otimes U$
			& $4\operatorname{Re}\!\bigl[(U_t)_{rr}\,
			\overline{(U_0)_{rr}}\bigr]$
			& Yes & $\tr(I_{4}) = 4$\\
			C & $I_{2}\otimes U\otimes I_{2}$
			& $4\operatorname{Re}\!\bigl[(U_t)_{rr}\,
			\overline{(U_0)_{rr}}\bigr]$
			& Yes & $\tr(I_{2})^{2} = 4$\\
			\bottomrule
		\end{tabular}
	\end{table}
	
	% ============================================================
	\section{The Two-Observable Decomposition Structure}
	\label{sec:two-obs-structure}
	% ============================================================
	
	\subsection{Setup and Definitions}
	\label{sec:setup}
	
	Theorem~\ref{thm:two-obs-completeness} (Sec.~\ref{sec:theorem-completeness})
	establishes that the full gauge-link correlator decomposes exactly into
	two independent bilinear contributions.  The defining equations are
	recalled below; per Q1 the original (6.1a)--(6.1c) sub-letter triplet
	has been split into three independently numbered equations:
	
	\begin{align}
		\Gdiag(t)  &\;=\; 8\!\sum_{\mathrm{links}}\,
		\operatorname{Re}[a_t\,\bar{a}_0],
		\label{eq:Gdiag-def-recall}\\[2pt]
		\Gcross(t) &\;=\; 8\!\sum_{\mathrm{links}}\,
		\operatorname{Re}[b_t\,\bar{b}_0],
		\label{eq:Gcross-def-recall}\\[2pt]
		\Gfull(t)  &\;=\; \Gdiag(t) + \Gcross(t).
		\label{eq:two-obs-completeness-recall}
	\end{align}
	Equation~\eqref{eq:two-obs-completeness-recall} is exact for every
	$\su(2)$ gauge configuration and every timeslice~$t$
	(Theorem~\ref{thm:two-obs-completeness}).
	
	% ------------------------------------------------------------
	\subsection{Two-Observable Decomposition}
	\label{sec:cor-two-obs-decomp}
	% ------------------------------------------------------------
	
	\begin{corollary}[Two-Observable Decomposition in the
		$\cl(0,2)\cong\HH$ Framework]
		\label{cor:two-obs-decomp}
		The following holds exactly for every $\su(2)$ field configuration
		and every~$t$:
		\begin{enumerate}
			\item[(i)] \textbf{Exactness}
			\textsc{[new --- proved as
				Theorem~\ref{thm:two-obs-completeness}]}:
			$\;\Gdiag(t) + \Gcross(t) = \Gfull(t).$
			
			\item[(ii)] \textbf{$a$-Channel Faithfulness}
			\textsc{[recalled --- Proposition~1 of
				Paper~4~\cite{FirestoneP4}, \S4.2]}:
			$\Gdiag$ depends only on $a_t$ and $\bar{a}_0$.
			
			\item[(iii)] \textbf{$b$-Channel Faithfulness}
			\textsc{[new --- proved as
				Proposition~\ref{prop:off-diag-collapse}]}:
			$\Gcross$ depends only on $b_t$ and $\bar{b}_0$.
		\end{enumerate}
		Corollary~\ref{cor:two-obs-decomp} follows immediately from
		Theorem~\ref{thm:two-obs-completeness},
		Proposition~1 of Paper~4~\cite{FirestoneP4} (\S4.2), and
		Proposition~\ref{prop:off-diag-collapse}.
		Parts~(i) and~(iii) are new results established in this paper;
		part~(ii) assembles Proposition~1 of Paper~4 for completeness of
		statement.
	\end{corollary}
	
	\begin{remark}[Parametric completeness]
		\label{rem:parametric-completeness}
		The pair $\{a,b\}$ with $|a|^{2}+|b|^{2}=1$ parametrises
		$\su(2)\cong S^{3}$.  This is a conjunction of~(ii), (iii), and
		the unitarity constraint $|a|^{2}+|b|^{2}=1$
		(labelled~(E4) in Paper~4~\cite{FirestoneP4}, retained here
		for cross-paper consistency)
		---not an independent result---but it justifies describing the
		decomposition as complete.
	\end{remark}
	
	\begin{remark}[Weak-coupling channel identification]
		\label{rem:weak-coupling-id}
		At leading order in $g$:
		$a = 1 + (ig/2)\,A^{3} + O(g^{2})$ and
		$b = (ig/\sqrt{2})\,W^{+} + O(g^{2})$
		(Eqs.~\eqref{eq:a-expansion},~\eqref{eq:b-expansion-W}).
		Therefore $\Gdiag$ probes $A^{3}$ and $\Gcross$ probes $W^{\pm}$.
		This identification is perturbative---valid to $O(g^{2})$ in the
		bilinears---and does not constitute part of the exact
		decomposition of Corollary~\ref{cor:two-obs-decomp}.
	\end{remark}
	
	\begin{proof}[Proof of Corollary~\ref{cor:two-obs-decomp}]
		Immediate from Theorem~\ref{thm:two-obs-completeness}
		(Sec.~\ref{sec:theorem-completeness}), Proposition~1 of
		Paper~4~\cite{FirestoneP4} (\S4.2), and
		Proposition~\ref{prop:off-diag-collapse}
		(Sec.~\ref{sec:off-diag-collapse}).
	\end{proof}
	
	% ------------------------------------------------------------
	\subsection{Why \texorpdfstring{$\{\Gdiag,\Gcross\}$}{\{Gdiag, Gcross\}}
		and Not \texorpdfstring{$\{G_{+},G_{-}\}$}{\{G+, G-\}}}
	\label{sec:why-G}
	% ------------------------------------------------------------
	
	By Corollary~1 of Paper~4~\cite{FirestoneP4}, \S4.4, $G_{+}(t) =
	G_{-}(t)$ identically for every $U\in\su(2)$.  Their sum is
	$G_{+} + G_{-} = 2 G_{+} = \Gdiag$, and the pair $\{G_{+},G_{-}\}$
	carries exactly one independent observable---the $a$-channel
	bilinear---and no information about $b$.  The apparent
	two-dimensional structure collapses to a single observable by
	Corollary~1 of Paper~4, leaving the entire $b$-channel inaccessible.
	
	The pair $\{\Gdiag,\Gcross\}$, by contrast, provides two genuinely
	independent observables: $\Gdiag$ is faithful to $a$ (part~(ii)) and
	$\Gcross$ is faithful to $b$ (part~(iii)), with the two channels
	linked only by the unitarity constraint $|a|^{2}+|b|^{2}=1$, which
	imposes no algebraic relation between $\operatorname{Re}[a_t\,
	\bar{a}_0]$ and $\operatorname{Re}[b_t\,\bar{b}_0]$ across timeslices.
	
	% ------------------------------------------------------------
	\subsection{Sign Non-Definiteness}
	\label{sec:prop-sign-nondef}
	% ------------------------------------------------------------
	
	\begin{proposition}[Sign Non-Definiteness]
		\label{prop:sign-nondef}
		$\Gdiag(t)$ and $\Gcross(t)$ are individually not sign-definite
		per configuration:
		\begin{enumerate}
			\item[(a)] $\Gdiag < 0$ is possible.  Example:
			$a_t = -1/\sqrt{2} + i\varepsilon$, $a_0 = 1/\sqrt{2}$;
			then $\operatorname{Re}[a_t\,\bar{a}_0] = -1/2 < 0$.
			
			\item[(b)] $\Gcross < 0$ is possible.  Example:
			$b_t = -1/\sqrt{2}$, $b_0 = 1/\sqrt{2}$;
			then $\operatorname{Re}[b_t\,\bar{b}_0] = -1/2 < 0$.
			
			\item[(c)] $\Gfull$ is not constrained positive per link by
			unitarity alone.
		\end{enumerate}
	\end{proposition}
	
	\begin{proof}
		Parts~(a) and~(b) follow by direct substitution; both examples
		satisfy $|a|^{2}+|b|^{2}=1$.  Part~(c) follows from~(a), (b), and
		Eq.~\eqref{eq:two-obs-completeness-recall}.  The numerical
		observation $\Dcross/\Gfull > 1$ for some $(\beta,t)$
		pairs~\cite{FirestoneP4} implies $\Gdiag < 0$ in those cases.
		$\Gcross < 0$ was observed in 37--62\% of (configuration,~$t$)
		pairs~\cite{FirestoneP4}.
	\end{proof}
	
	% ------------------------------------------------------------
	\subsection{Unitarity Constraint and Functional Independence}
	\label{sec:cor-unitarity}
	% ------------------------------------------------------------
	
	\begin{corollary}[Unitarity Constraint]
		\label{cor:unitarity}
		The $\su(2)$ constraint $|a|^{2}+|b|^{2}=1$ holds per link per
		timeslice.  It imposes no algebraic relation between
		$\operatorname{Re}[a_t\,\bar{a}_0]$ and
		$\operatorname{Re}[b_t\,\bar{b}_0]$ across distinct timeslices
		$t \neq 0$.  The two channels are functionally independent
		observables.
	\end{corollary}
	
	\begin{proof}
		The constraint $|a_t|^{2}+|b_t|^{2}=1$ bounds moduli individually
		but places no constraint on the phases $\arg(a_t\,\bar{a}_0)$ or
		$\arg(b_t\,\bar{b}_0)$, which determine the real parts of the
		bilinears.  For fixed $t=0$ reference link and varying $t > 0$,
		the values of $\operatorname{Re}[a_t\,\bar{a}_0]$ and
		$\operatorname{Re}[b_t\,\bar{b}_0]$ range independently over
		$[-1,1]$.
	\end{proof}
	
	% ============================================================
	\section{Connection to \texorpdfstring{$\su(2)_L$}{SU(2)\_L} /
		\texorpdfstring{$\su(2)_R$}{SU(2)\_R}}
	\label{sec:SU2-LR}
	% ============================================================
	
	\subsection{Channel Assignment at \texorpdfstring{$O(g)$}{O(g)}}
	\label{sec:channel-assign}
	
	The three $\su(2)$ generators in the fundamental representation are
	$T^{3} = \mathrm{diag}(1,-1)/2$, $T_{+} = \sigma_{+}/2$,
	$T_{-} = \sigma_{-}/2$.  The key point is that $T^{3}$ is purely
	diagonal and $T_{\pm}$ are purely off-diagonal.  Consequently, at
	$O(g)$ in $U = I + i g\,a_{\mathrm{lat}}\,A^{c}_{\mu} T^{c} + O(g^{2})$:
	
	\begin{itemize}
		\item $g\,A^{3}_{\mu}\,T^{3}$ contributes only to $(U)_{00}$ and
		$(U)_{11}$ ($a$-channel only).
		
		\item $g\,W^{+}_{\mu}\,T_{-}$ and $g\,W^{-}_{\mu}\,T_{+}$
		contribute only to $(U)_{01}$ and $(U)_{10}$ ($b$-channel only).
	\end{itemize}
	
	This separation is exact at $O(g)$.  Checks D17-R1, D17-R2, D17-R3
	verify: R1 confirms
	$\operatorname{Im}(a_\varepsilon - 1)/\varepsilon \to \alpha^{3}/2$
	(only $T^{3} A^{3}$); R2 confirms
	$\partial(\operatorname{Im} a)/\partial \alpha^{1} = 0$
	(no $A^{1}$ contamination); R3 confirms
	$\partial(\operatorname{Im} a)/\partial \alpha^{2} = 0$
	(no $A^{2}$ contamination).
	
	\begin{table}[ht]
		\centering
		\caption{Leading-order $\su(2)$ channel assignment.  Each
			generator--field combination contributes to exactly one
			matrix element of $U$ at $O(g)$, mapping cleanly onto either
			the $a$- or $b$-channel.}
		\label{tab:Og-channel-assignment}
		\small
		\setlength{\tabcolsep}{8pt}
		\begin{tabular}{@{}llll@{}}
			\toprule
			\textbf{Coupling} & \textbf{Matrix element}
			& \textbf{Channel} & \textbf{$O(g)$ contribution}\\
			\midrule
			$T^{3}\,A^{3}_{\mu}$ & $(U)_{00} = a$
			& $a$-channel & $+(ig/2)\,A^{3}$\\
			$T_{-}\,W^{+}_{\mu}$ & $(U)_{01} = b$
			& $b$-channel & $+(ig/\sqrt{2})\,W^{+}$\\
			$T_{+}\,W^{-}_{\mu}$ & $(U)_{10} = -b^{*}$
			& $b$-channel & $+(ig/\sqrt{2})\,W^{-}$\\
			$T^{3}\,A^{3}_{\mu}$ & $(U)_{11} = a^{*}$
			& $a$-channel & $-(ig/2)\,A^{3}$\\
			\bottomrule
		\end{tabular}
	\end{table}
	
	% ------------------------------------------------------------
	\subsection{The \texorpdfstring{$b$}{b}-Channel as Necessary Conduit}
	\label{sec:b-conduit}
	% ------------------------------------------------------------
	
	\begin{proposition}[$b$-Channel Necessary Conduit]
		\label{prop:b-conduit}
		\begin{itemize}
			\item[(b1)]
			\textbf{Necessity.}  Any gauge-link correlator sensitive to
			$W^{\pm}$ at $O(g)$ must involve the $b$-channel.  No diagonal
			projector correlator of the form
			$\operatorname{Re}\,\tr[P_{r}\, E_{t}\, P_{r}\, E_{0}^{\dagger}]$
			can access $b$ (Lemma~\ref{lem:b-blindness}).
			
			\item[(b2)]
			\textbf{Non-sufficiency.}  The $b$-channel does not
			distinguish $\su(2)_L$ from $\su(2)_R$.  The link variable
			$U$ transforms identically under $\su(2)$ regardless of the
			handedness of any matter field.  Chirality is encoded in the
			matter sector only.
		\end{itemize}
	\end{proposition}
	
	\begin{proof}
		\emph{Proof of (b1).}  By Lemma~\ref{lem:b-blindness}, the
		diagonal projector correlator is independent of $b$ for every
		embedding case.  Since $W^{\pm}$ enter only through $b$ at $O(g)$
		(Table~\ref{tab:Og-channel-assignment}), no diagonal projector
		correlator can be sensitive to $W^{\pm}$.
		$\Gcross = 8\sum\operatorname{Re}[b_t\,\bar{b}_0]$ is sensitive
		to $W^{\pm}$ by~\eqref{eq:b-expansion-W} and is the unique
		algebraic structure accessing $b$
		(Proposition~\ref{prop:off-diag-collapse}).
		
		\emph{Proof of (b2).}  The pure-gauge link variable carries no
		record of fermion handedness; this is determined by the
		transformation properties of fermion fields under Lorentz and
		gauge symmetry.
	\end{proof}
	
	\medskip
	\noindent
	\emph{The embedding does not, by itself, distinguish $\su(2)_L$ from
		$\su(2)_R$; this requires external fermionic input.}
	
	% ------------------------------------------------------------
	\subsection{Conditional Prediction for Fermionic Observables
		--- Programme Target}
	\label{sec:fermion-target}
	% ------------------------------------------------------------
	
	\noindent
	\textbf{This subsection describes a programme target for future work.
		It is not a result of Paper~5.  Paper~5 contains no fermion action
		and makes no claim of testing the following prediction.}
	
	\medskip
	Suppose a fermionic propagator $S(x,y)$ is introduced on the same
	lattice, transforming in the fundamental representation of $\su(2)$.
	Define the cross-projected propagator:
	\begin{equation}
		S_{\mathrm{cross}}(x,y) \;=\; \PLW_+ \, S(x,y) \, \PLW_-.
		\label{eq:Scross-def}
	\end{equation}
	The conditional prediction of the $b$-channel framework is:
	
	\begin{itemize}
		\item For an $\su(2)_L$ doublet fermion: $S_{\mathrm{cross}}
		\neq 0$ (the doublet mixes isospin components via the
		$T_{-}\,W^{+}_{\mu}$ coupling).
		
		\item For an $\su(2)_R$ singlet fermion: $S_{\mathrm{cross}} = 0$
		(a singlet carries no $\su(2)$ index).
	\end{itemize}
	
	The Nielsen--Ninomiya obstruction~\cite{NielsenNinomiya1981a,%
		NielsenNinomiya1981b} states that no local, Hermitian,
	translation-invariant lattice fermion action can simultaneously be
	free of fermion doublers and chirally symmetric.  Implementation
	requires domain wall~\cite{Kaplan1992} or overlap
	fermions~\cite{Neuberger1998} (Ginsparg--Wilson).
	This is designated Future Requirements F1--F4 and constitutes the
	entry point for Paper~6.
	
	% ============================================================
	\section{Gauge Properties of \texorpdfstring{$\Dcross$}{Delta\_cross}
		and the Landau Gauge Prescription}
	\label{sec:gauge-properties}
	% ============================================================
	
	% NOTE: Subsections §8a-§8d in the original draft used letter-style
	% labels. These have been re-rendered as §8.1-§8.4 for consistency
	% with the numeric subsection style used in all other sections.
	
	\subsection{Proof of Gauge Non-Invariance}
	\label{sec:gauge-noninv}
	
	\begin{lemma}[Transformation of $b$ under Gauge Rotations]
		\label{lem:b-transform}
		Under $U' = g\, U\, h^{\dagger}$, with
		$g = \begin{pmatrix}\alpha & \beta\\ -\beta^{*} & \alpha^{*}
		\end{pmatrix}$ and
		$h = \begin{pmatrix}\gamma & \delta\\ -\delta^{*} & \gamma^{*}
		\end{pmatrix}$ elements of $\su(2)$, the transformed off-diagonal
		element is
		\begin{equation}
			b' = (U')_{01}
			\;=\; \alpha\gamma\,b + \beta\gamma\,a^{*}
			- \alpha\delta\,a + \beta\delta\,\bar{b}.
			\label{eq:b-transform}
		\end{equation}
	\end{lemma}
	
	\begin{proof}
		Direct matrix multiplication of
		$(g\, U\, h^{\dagger})_{01}
		= \sum_{j,k}(g)_{0j}\,(U)_{jk}\,(h^{\dagger})_{k1}$,
		substituting the $\su(2)$ parametrisations.
		\emph{Counterexample:}
		$g_t = i\sigma_{2}
		= \begin{pmatrix}0 & 1\\ -1 & 0\end{pmatrix}$,
		$h_t = I_2$ gives $b'_t = a^{*}_t$, demonstrating that $b$ is
		not gauge-invariant.  Verified on 50 random $\su(2)$ pairs by
		check D18b-f (residual $0.00\times 10^{+00}$).
	\end{proof}
	
	\begin{proposition}[$\Dcross$ is Not Gauge-Invariant]
		\label{prop:gauge-noninv}
		Under the lattice gauge transformation
		$U_\mu(x) \to g(x)\,U_\mu(x)\,g^{\dagger}(x+\hat{\mu})$, with
		$g(x)\in\su(2)$, $\Dcross(t) = 8\sum
		\operatorname{Re}[b_t\,\bar{b}_0]$ changes.
	\end{proposition}
	
	\begin{proof}
		Under the gauge transformation with $g_t \neq g_0$, the
		transformed elements $b'_t$ and $b'_0$ each depend on $(g_t,g_0)$
		through independent linear combinations of
		$a_t,\,b_t,\,a^{*}_t,\,\bar{b}_t$ and
		$a_0,\,b_0,\,a^{*}_0,\,\bar{b}_0$ respectively
		(Lemma~\ref{lem:b-transform}).
		$\operatorname{Re}[b'_t\,\bar{b}'_0] \neq
		\operatorname{Re}[b_t\,\bar{b}_0]$ generically.
		\emph{Counterexample:} $g_t = i\sigma_{2}$, $g_0 = I_2$ gives
		$b'_t = a^{*}_t$, $b'_0 = b_0$, so
		$\operatorname{Re}[b'_t\,\bar{b}'_0]
		= \operatorname{Re}[a^{*}_t\,\bar{b}_0]
		\neq \operatorname{Re}[b_t\,\bar{b}_0]$ generically.
	\end{proof}
	
	Check D18a (inverted) verifies gauge non-invariance on 50
	configurations using Haar-random gauge transformations, finding
	min residual $1.07\times 10^{-2}\gg 10^{-6}$.  Check D18b-n verifies
	the specific counterexample.  Scope: D18 is designed for Haar-random
	transforms; near-identity transforms
	($g \approx I + i\varepsilon H$, $\varepsilon \lesssim 10^{-7}$) not
	tested.  Proposition~\ref{prop:gauge-noninv} proves non-invariance
	analytically for all transforms without restriction on the magnitude
	of~$g$.
	
	% ------------------------------------------------------------
	\subsection{Landau Gauge Prescription}
	\label{sec:landau-presc}
	% ------------------------------------------------------------
	
	Since $\Dcross$ is gauge-dependent
	(Proposition~\ref{prop:gauge-noninv}), a gauge-fixing prescription
	is required.  The standard choice for gluon propagator studies is
	Landau gauge, defined by the maximisation condition:
	\begin{equation}
		F_{L}[\{g(x)\}]
		\;=\; \operatorname{Re}\sum_{x,\mu}
		\tr[U^{g}_{\mu}(x)]
		\;\longrightarrow\; \text{maximum over } \{g(x)\},
		\label{eq:Landau-gauge}
	\end{equation}
	where $U^{g}_{\mu}(x) = g(x)\,U_{\mu}(x)\,g^{\dagger}(x+\hat{\mu})$.
	The Landau condition $\partial_{\mu} A^{a}_{\mu} = 0$ is recovered in
	the continuum limit.  The target comparison data are the $\su(2)$
	Landau-gauge gluon propagator computations of Cucchieri and
	Mendes~\cite{CucchieriMendes2008} and Bogolubsky et
	al.~\cite{Bogolubsky}.
	
	% ------------------------------------------------------------
	\subsection{Formal Comparison to the
		\texorpdfstring{$W^{\pm}$}{W±} Propagator}
	\label{sec:Wpm-compare}
	% ------------------------------------------------------------
	
	In Landau gauge, the $W^{\pm}$ propagator at zero spatial momentum
	is:
	\[
	C_{W^{\pm}}(t)
	\;=\; \frac{1}{V}\sum_{x}
	\bigl\langle W^{+}_{i}(x,t)\,W^{-}_{i}(x,0)\bigr\rangle.
	\]
	Using the weak-coupling identification~\eqref{eq:b-expansion-W},
	$b \approx (ig/\sqrt{2})\,W^{+}$, so
	$\bar{b}_{0} \approx (-ig/\sqrt{2})\,W^{-}_{0}$, giving
	\[
	b_{t}\,\bar{b}_{0}
	\;\approx\; (ig/\sqrt{2})\,(-ig/\sqrt{2})\,W^{+}_{t}\,W^{-}_{0}
	\;=\; +(g^{2}/2)\,W^{+}_{t}\,W^{-}_{0}.
	\]
	The sign is \textbf{positive: $+(g^{2}/2)$}, confirmed
	by~\eqref{eq:b-bilinear-Og2}:
	$\operatorname{Re}[b_t\,\bar{b}_0]
	= (g^2/4)\,(A^{1}_t A^{1}_0 + A^{2}_t A^{2}_0) \geq 0$ at $O(g^2)$.
	The formal identification in Landau gauge is:
	\begin{equation}
		\Dcross(t)\big|_{\mathrm{Landau}}\big/\,8V
		\;\approx\; +(g^{2}/2)\cdot C_{W^{\pm}}(t) + O(g^{3}).
		\label{eq:Dcross-Landau}
	\end{equation}
	
	Four caveats apply:
	\begin{itemize}
		\item[(i)]
		\textbf{Perturbative validity.}
		Equation~\eqref{eq:Dcross-Landau} holds at $O(g^{2})$ and does
		not account for non-perturbative contributions.  At
		$\beta_{\mathrm{eff}}\in\{5.0,\,6.0,\,9.0,\,12.0\}$, the
		$O(g^{2})$ approximation is not quantitatively valid; the $O(1)$
		ratio $\Dcross/\Gfull\in[0.72,\,2.33]$ evidences higher-order
		contributions~\cite{FirestoneP4}.
		
		\item[(ii)]
		\textbf{Spatial links only.}
		Papers~4 and~5 use temporal gauge $A_{0}=0$, fixing temporal
		links to $I_{2}$.  Full Landau gauge requires temporal links;
		Eq.~\eqref{eq:Dcross-Landau} applies to spatial links only in
		the current implementation.
		
		\item[(iii)]
		\textbf{Strong-coupling breakdown.}
		The $O(g^{2})$ approximation requires lattice renormalisation
		to be applied systematically at finite $\beta_{\mathrm{eff}}$.
		
		\item[(iv)]
		\textbf{Gribov copy dependence --- systematic uncertainty.}
		The Landau gauge functional $F_{L}$ generically possesses
		multiple local maxima on the gauge orbit (Gribov copies).  Each
		copy gives a different value of $\Dcross$.  This systematic
		uncertainty is distinct from caveat~(i) and is not addressed in
		Paper~5.  This is the primary open question OQ1
		(Sec.~\ref{sec:discussion}).  See Sec.~\ref{sec:future-sim} for
		the proposed handling strategy.
	\end{itemize}
	
	% ------------------------------------------------------------
	\subsection{Future Simulation Requirements}
	\label{sec:future-sim}
	% ------------------------------------------------------------
	
	The following requirements are designated for Paper~6 and are not
	implemented in Paper~5:
	
	\begin{description}
		\item[F5.] Steepest-descent or Fourier-accelerated gauge fixing
		to implement~\eqref{eq:Landau-gauge} on existing ensembles.
		
		\item[F6.] Fourier-accelerated gauge fixing as an alternative for
		improved convergence at large volumes.
		
		\item[F7.] Gribov copy strategy: $\geq 5$ independent random
		gauge starts per configuration; select copy achieving maximum
		$F_{L}$; flag configurations where $F_{L}$ variation across
		copies exceeds $1\%$.
		
		\item[F8.] Post-gauge-fix measurement of $\Dcross$ on
		gauge-fixed ensembles.
		
		\item[F9.] Propagator extraction and quantitative comparison to
		the $\su(2)$ Landau-gauge gluon propagator data
		of~\cite{CucchieriMendes2008,Bogolubsky}.
	\end{description}
	
	% ============================================================
	\section{Discussion}
	\label{sec:discussion}
	% ============================================================
	
	This section synthesises the results of
	Sections~\ref{sec:framework}--\ref{sec:gauge-properties} in prose,
	draws five conclusions, and identifies seven open questions.  No new
	algebra or numerical results are introduced.
	
	% ------------------------------------------------------------
	\subsection{C1: Exactness and the Algebraic Status of
		\texorpdfstring{$\Gcross$}{Gcross}}
	\label{sec:disc-exactness}
	% ------------------------------------------------------------
	
	The principal achievement of this paper is the exact algebraic
	identification of $\Gcross$.  Prior to
	Proposition~\ref{prop:off-diag-collapse}, the cross term $\Dcross$
	was defined subtractionally in Paper~4~\cite{FirestoneP4}, and its
	$O(1)$ magnitude ($\Dcross/\Gfull\in[0.72,\,2.33]$~\cite{FirestoneP4})
	was observed without an algebraic explanation.
	Theorem~\ref{thm:two-obs-completeness} provides that explanation:
	$\Gcross = 8\sum\operatorname{Re}[b_t\,\bar{b}_0]$ exactly, for every
	$\su(2)$ configuration and every timeslice, where $b = (U)_{01}$ is
	the off-diagonal matrix element.
	
	Corollary~\ref{cor:two-obs-decomp} assembles this with the
	$a$-channel faithfulness result of Proposition~1 of
	Paper~4~\cite{FirestoneP4} into a complete statement of the exact
	algebraic decomposition.  The word ``complete'' in
	Remark~\ref{rem:parametric-completeness} of
	Corollary~\ref{cor:two-obs-decomp} carries a precise meaning:
	together, $\Gdiag$ and $\Gcross$ are sensitive to every degree of
	freedom in $U\in\su(2)\cong S^{3}$.  This algebraic completeness does
	not imply gauge invariance, sign-definiteness, or physical
	observability without gauge fixing.
	
	% ------------------------------------------------------------
	\subsection{C2: Embedding Independence}
	\label{sec:disc-embedding}
	% ------------------------------------------------------------
	
	Proposition~\ref{prop:rank1-univ} establishes that the rank-1
	collapse $\Pir\,U\,\Pir = (U)_{rr}\,\Pir$ holds identically for all
	three natural tensor-product embeddings of $U$ into $\Cplx^{8}$
	(Cases~A, B, C).  The spectator dimensions contribute a universal
	trace factor of~$4$, and the extracted bilinear is the same function
	of $(U_{t})_{rr}$ regardless of embedding.  The two-observable
	structure of $\{\Gdiag,\Gcross\}$ is intrinsic to the fundamental
	spin-$1/2$ representation of $\su(2)$, confirmed at 1~ULP residual
	by checks D14$^{\pm}$, D15$^{\pm}$.
	
	% ------------------------------------------------------------
	\subsection{C3: Physical Channel Assignment}
	\label{sec:disc-channel}
	% ------------------------------------------------------------
	
	At $O(g)$ in the weak-coupling expansion, the channel assignment is
	exact by group-theoretic argument: $T^{3}$ is purely diagonal and
	contributes solely to the $a$-channel; $T_{\pm}$ are purely
	off-diagonal and contribute solely to the $b$-channel.  There is no
	cross-contamination at this order, confirmed by checks D17-R1,
	D17-R2, D17-R3.  The non-perturbative contributions visible in the
	simulated ensembles are consistent with this assignment at the
	algebraic level, though quantitative identification requires Landau
	gauge fixing and is subject to the caveats of
	Sec.~\ref{sec:Wpm-compare}.
	
	% ------------------------------------------------------------
	\subsection{C4: Gauge Properties and the Path to a Physical
		Observable}
	\label{sec:disc-gauge}
	% ------------------------------------------------------------
	
	The gauge non-invariance of $\Gcross$ proved in
	Proposition~\ref{prop:gauge-noninv} is expected: charged-current
	correlators are not gauge-invariant in pure Yang--Mills theory
	without gauge-invariant composite operators or a physical gauge
	condition.  This is a standard feature, not a pathology of the
	framework.
	
	The path to a physical observable is through Landau gauge fixing
	(Secs.~\ref{sec:landau-presc}--\ref{sec:future-sim}).  In Landau
	gauge, the formal identification
	$\Dcross(t)|_{\mathrm{Landau}}/8V \approx +(g^2/2)\,C_{W^{\pm}}(t)$
	at $O(g^{2})$ connects $\Gcross$ to the zero-spatial-momentum
	$W^{\pm}$ temporal correlator, with~\cite{CucchieriMendes2008,%
		Bogolubsky} as comparison targets.  The positive sign $+(g^{2}/2)$
	---confirmed by~\eqref{eq:b-bilinear-Og2} and
	Sec.~\ref{sec:Wpm-compare}---means $\Gcross$ is non-negative at
	leading perturbative order; the observed negative values on
	thermalized configurations (37--62\% of cases~\cite{FirestoneP4}) are
	a genuine non-perturbative effect.  Gribov copy dependence introduces
	a systematic uncertainty not bounded in this paper, designated OQ1.
	
	% ------------------------------------------------------------
	\subsection{C5: The \texorpdfstring{$b$}{b}-Channel as Necessary
		Conduit; Matter for Chirality}
	\label{sec:disc-chirality}
	% ------------------------------------------------------------
	
	Proposition~\ref{prop:b-conduit} establishes two complementary facts.
	First, the $b$-channel is an unavoidable algebraic gateway for any
	gauge-link correlator sensitive to $W^{\pm}$: no diagonal projector
	correlator can access $b$ (Lemma~\ref{lem:b-blindness}), so $\Gcross$
	is the unique algebraic structure through which the charged gauge
	boson sector is visible.  Second, the $b$-channel does not by itself
	select chirality: the pure-gauge link variable transforms identically
	under $\su(2)$ regardless of matter handedness, and the distinction
	between left- and right-handed couplings is encoded entirely in the
	matter sector.  The identification of $\Gcross$ as the algebraic
	conduit for $W^{\pm}$ is a necessary but not sufficient condition for
	implementing the chiral structure of the electroweak interaction.
	Chirality selection requires fermionic fields with definite
	handedness, and faces the Nielsen--Ninomiya
	obstruction~\cite{NielsenNinomiya1981a,NielsenNinomiya1981b} in any
	local, Hermitian, doubler-free formulation.
	
	% ------------------------------------------------------------
	\subsection{Open Questions OQ1--OQ7}
	\label{sec:disc-OQ}
	% ------------------------------------------------------------
	
	\begin{description}
		\item[OQ1] \emph{(Landau gauge measurement.)}
		Implementing requirements F5--F9 on the existing ensembles and
		performing a quantitative comparison to~\cite{%
			CucchieriMendes2008,Bogolubsky}.  The Gribov copy systematic
		(caveat~(iv) of Sec.~\ref{sec:Wpm-compare}) must be characterised
		before the comparison can be made quantitatively reliable.
		
		\item[OQ2] \emph{(Sign-definiteness of ensemble averages.)}
		Whether $\langle\Gdiag\rangle$ and $\langle\Gcross\rangle$ are
		sign-definite at finite $\beta$ is an open numerical question
		requiring measurements on larger ensembles.
		
		\item[OQ3] \emph{($\su(3)$ extension.)}
		Extending the two-observable decomposition to $\su(3)$ within
		the $\cl(6)$ framework of Paper~1~\cite{FirestoneP1}, where the
		algebra is larger and the projector structure more complex.
		
		\item[OQ4] \emph{(Chirality and fermionic content.)}
		The conditional prediction of Sec.~\ref{sec:fermion-target}
		requires a chirally symmetric lattice fermion action.  The
		Nielsen--Ninomiya
		obstruction~\cite{NielsenNinomiya1981a,NielsenNinomiya1981b}
		requires domain wall or overlap fermions.  Designated Future
		Requirements F1--F4.
		
		\item[OQ5] \emph{(Electroweak mixing, $\uu(1)_Y$, and $W/Z$ mass
			ratio.)}
		Including $\uu(1)_Y$ and the Higgs mechanism.  Outside the scope
		of the present series.
		
		\item[OQ6] \emph{(Higher-spin and reducible representations.)}
		Extension of the channel decomposition to higher-spin or
		reducible representations of $\su(2)$.
		
		\item[OQ7] \emph{(Temporal gauge and full spacetime structure.)}
		A formulation incorporating all four-dimensional spacetime links
		in Landau gauge for proper continuum limit comparison and
		momentum-space propagator extraction.
	\end{description}
	
	\medskip\noindent
	OQ1--OQ4 are the most pressing; OQ5--OQ7 are longer-range.
	
	% ============================================================
	\section{Conclusions}
	\label{sec:conclusions}
	% ============================================================
	
	We have presented a two-observable decomposition of the $\su(2)$
	temporal gauge-link correlator within the $\cl(0,2)\cong\HH$
	algebraic framework.  Three results stand out as the principal
	contributions of this paper.  We emphasise that the quaternionic
	parametrisation $U = a\cdot I_2 + b\cdot J$ and the resulting
	bilinear split are standard; the new contributions are the
	\emph{projector formulation} of the channel separation, the
	\emph{algebraic closure of the Paper~4 gap} (giving $\Dcross$ an
	independent generator), and the \emph{Landau-gauge propagator
	identification} at $O(g^2)$.
	
	% ------------------------------------------------------------
	\subsection*{Result 1 --- Exact Algebraic Decomposition
		(Cor.~\ref{cor:two-obs-decomp};
		Thm.~\ref{thm:two-obs-completeness})}
	\label{sec:result1}
	% ------------------------------------------------------------
	
	The pair $\{\Gdiag(t),\Gcross(t)\}$ constitutes an exact algebraic
	decomposition of $\Gfull(t)$ in the $\cl(0,2)\cong\HH$ framework.
	Specifically, $\Gdiag(t) + \Gcross(t) = \Gfull(t)$ holds exactly for
	every $\su(2)$ gauge configuration and every timeslice~$t$, with
	$\Gdiag(t) = 8\sum\operatorname{Re}[a_t\bar{a}_0]$ faithful to the
	$a$-channel (Proposition~1 of Paper~4~\cite{FirestoneP4}) and
	$\Gcross(t) = 8\sum\operatorname{Re}[b_t\bar{b}_0]$ faithful to the
	$b$-channel
	(Proposition~\ref{prop:off-diag-collapse}).  This resolves the
	structural gap of Paper~4: the cross term $\Dcross$, previously
	defined subtractionally, is exactly the $b$-channel bilinear.
	Verified to $|\Gdiag + \Gcross - \Gfull|\leq 1.78\times 10^{-15}$
	and $|\Gcross - \Dcross|\leq 1.14\times 10^{-13}$ across all four
	ensembles (D7a--D7c, D8\,$\times\,4\beta$).
	
	% ------------------------------------------------------------
	\subsection*{Result 2 --- Embedding Independence
		(Prop.~\ref{prop:rank1-univ})}
	\label{sec:result2}
	% ------------------------------------------------------------
	
	The rank-1 collapse $\Pir\, U\,\Pir = (U)_{rr}\,\Pir$ --- which
	forces the $b$-channel blindness of all diagonal projector
	correlators and thereby necessitates the off-diagonal structure of
	$\Gcross$ --- is intrinsic to the fundamental spin-$1/2$
	representation of $\su(2)$.  It holds identically for every
	tensor-product embedding of $U$ into $\Cplx^{8}$ (Cases~A, B, C of
	Sec.~\ref{sec:embedding-univ}), with the same bilinear formula and
	the same spectator trace factor of~$4$ in each case.  The
	two-observable structure of $\{\Gdiag,\Gcross\}$ is not an artifact
	of the specific embedding $I_{2}\otimes U\otimes I_{2}$.
	
	% ------------------------------------------------------------
	\subsection*{Result 3 --- Physical Channel Assignment and Gauge
		Status (Prop.~\ref{prop:b-conduit};
		Prop.~\ref{prop:gauge-noninv})}
	\label{sec:result3}
	% ------------------------------------------------------------
	
	At $O(g)$, the generator structure of $\su(2)$ enforces an exact
	channel separation: $T^{3}\,A^{3}_{\mu}$ contributes only to the
	$a$-channel and $T_{\pm}\,W^{\mp}_{\mu}$ contribute only to the
	$b$-channel (D17-R1, D17-R2, D17-R3).  $\Gdiag$ probes $A^{3}$;
	$\Gcross$ probes $W^{\pm}$.  The $b$-channel is a necessary algebraic
	conduit for all charged-current coupling
	(Proposition~\ref{prop:b-conduit}).  $\Gcross$ is proved not to be
	gauge-invariant (Proposition~\ref{prop:gauge-noninv}).  In Landau
	gauge:
	\[
	\Dcross(t)\big|_{\mathrm{Landau}}\big/\,8V
	\;\approx\; +(g^{2}/2)\cdot C_{W^{\pm}}(t) + O(g^{3}),
	\]
	with positive sign $+(g^{2}/2)$ confirmed by the weak-coupling
	calculation of Sec.~\ref{sec:Wpm-compare}.
	
	% ------------------------------------------------------------
	\subsection*{What Paper~5 Does Not Establish}
	\label{sec:not-establish}
	% ------------------------------------------------------------
	
	Paper~5 does not establish gauge invariance of $\Gcross$ or
	$\Dcross$.  It does not implement Landau gauge fixing, does not
	compute the $W^{\pm}$ propagator numerically, and does not perform a
	quantitative comparison to the lattice data
	of~\cite{CucchieriMendes2008,Bogolubsky}.  It does not address
	electroweak mixing, $W/Z$ mass ratios, the Higgs mechanism, or
	chiral symmetry breaking.  It does not implement a fermionic action
	of any kind and makes no claim of overcoming the Nielsen--Ninomiya
	obstruction~\cite{NielsenNinomiya1981a,NielsenNinomiya1981b}.
	\emph{The embedding does not, by itself, distinguish $\su(2)_L$ from
		$\su(2)_R$; this requires external fermionic input.}
	All 51 algebraic and numerical checks (D1--D18) are satisfied at
	tolerance $\leq 10^{-13}$ or better; the results are exact in the
	algebraic sense of Corollary~\ref{cor:two-obs-decomp} and
	Theorem~\ref{thm:two-obs-completeness}, perturbative in the physical
	identification of Sec.~\ref{sec:Wpm-compare}, and conditional in the
	programme targets of Sec.~\ref{sec:fermion-target}.
	
	% ============================================================
	%  APPENDICES
	% ============================================================
	\appendix
	\renewcommand{\thetable}{\Alph{section}.\arabic{table}}
	\numberwithin{table}{section}
	\renewcommand{\theequation}{\Alph{section}.\arabic{equation}}
	\numberwithin{equation}{section}
	
	% ------------------------------------------------------------
	\section{Algebraic Verification Checks}
	\label{app:A}
	% ------------------------------------------------------------
	
	\setcounter{table}{-1}   % first \begin{table} in this appendix --> Table A.0
		
		All 51 verification instances were executed across check IDs
		D1--D18.  The programme achieves a complete \textsc{pass} record:
		51/51 instances pass, 0 fail, with 1 flagged anomaly (D10b, closed)
		and 3 inverted checks (D9a, D18a, D18b-n).  The next available check
		identifier for Paper~6 is D19.
		
		\begin{table}[ht]
			\centering
			\caption{Complete check programme --- D1 through D18.}
			\label{tab:check-programme}
			\footnotesize
			\setlength{\tabcolsep}{4pt}
			\renewcommand{\arraystretch}{1.15}
			\begin{tabular}{@{}l c p{4.4cm} p{3.2cm} l l@{}}
				\toprule
				\textbf{Check ID} & \textbf{\S} & \textbf{Description}
				& \textbf{Expected} & \textbf{Tol.} & \textbf{Result}\\
				\midrule
				D1 (D1a--D1d) & \S\ref{sec:intro}--\S\ref{sec:framework}
				& Data structure; $\su(2)$ param; $|a|^2 + |b|^2 = 1$
				& Param exact; (E4) holds
				& $\leq 10^{-15}$ & \textsc{pass}\\
				D2 (D2a--D2d) & \S\ref{sec:framework}
				& Clifford grade decomp; $\cl^{0}\cong\Cplx$;
				$\cl^{1} = \Cjhat$; RM1--RM5
				& Grade split consistent
				& $\leq 10^{-15}$ & \textsc{pass}\\
				D3 (D3a--D3d)
				& \S\ref{sec:framework}--\S\ref{sec:two-obs-decomp-section}
				& Rank-1 collapse (diag.\ \& off-diag.);
				$b$-blindness of $P_r$
				& Collapse formulae exact
				& $\leq 10^{-15}$ & \textsc{pass}\\
				D4 (D4a--D4h) & \S\ref{sec:framework}
				& Generator commutation;
				Casimir $= (3/4)\,I_2$; algebra closure
				& Relations exact
				& $\leq 10^{-14}$ & \textsc{pass}\\
				D5 & \S\ref{sec:two-obs-decomp-section}
				& Prop.~\ref{prop:off-diag-collapse}:
				$\operatorname{Re}\tr[\PLW_+ E_t \PLW_- E_0^{\dagger}]
				= 4\operatorname{Re}[b_t \bar{b}_0]$, 50 pairs
				& RHS matches LHS
				& $\leq 10^{-14}$ & \textsc{pass}\\
				D6 & \S\ref{sec:two-obs-decomp-section}
				& Cor.~\ref{cor:orientation-sym}:
				orientation symmetry, 50 pairs
				& Equality holds
				& $\leq 10^{-14}$ & \textsc{pass}\\
				D7 (D7a--D7c) & \S\ref{sec:two-obs-decomp-section}
				& Thm.~\ref{thm:two-obs-completeness}:
				$|\Gdiag + \Gcross - \Gfull|$
				& $\leq 1.78\times 10^{-15}$
				& $\leq 10^{-14}$ & \textsc{pass}\\
				D8 ($\times\,4\beta$) & \S\ref{sec:two-obs-decomp-section}
				& $\Gcross = \Dcross$:
				$|\Gcross - \Dcross|$, all ensembles
				& $\leq 1.14\times 10^{-13}$
				& $\leq 10^{-12}$ & \textsc{pass}\\
				D9 (D9a) & \S\ref{sec:two-obs-decomp-section}
				& Sign non-definiteness of $\Gcross$ (inverted)
				& $\Gcross < 0$ in $\geq 1$ case
				& Inverted & \textsc{pass}\\
				D10 (D10a$\times 3\varepsilon$; D10b*)
				& \S\ref{sec:physical-interp}
				& $b$-channel WC error vs $\varepsilon$; log-log slope
				& $O(\varepsilon^{3})$; slope$=3.000$
				& $\leq 10^{-10}$ & \textsc{pass}*\\
				D11 (D11a$\times 3\varepsilon$; D11b)
				& \S\ref{sec:physical-interp}
				& $a$-channel WC error vs $\varepsilon$; log-log slope
				& $O(\varepsilon^{2})$; slope$=2.000$
				& $\leq 10^{-10}$ & \textsc{pass}\\
				D12 (D12b) & \S\ref{sec:physical-interp}
				& Monotone $|R(\beta)|$ at $t=5$ across ensembles
				& Monotone in $\beta_{\mathrm{eff}}$
				& Qualitat.\ & \textsc{pass}\\
				D13 (D13a) & \S\ref{sec:physical-interp}
				& Cold-start $\langle|b|^2\rangle\approx 0.5$ (Haar)
				& $\approx 0.5$
				& $\leq 10^{-2}$ & \textsc{pass}\\
				D14 (D14$^{+}$, D14$^{-}$) & \S\ref{sec:embedding-univ}
				& Case A ($U\otimes I_4$): $8\times 8$ trace vs
				$2\times 2$ element, $r=\pm$, 50 pairs
				& $\leq 2.22\times 10^{-16}$ (1 ULP)
				& 1 ULP & \textsc{pass}\\
				D15 (D15$^{+}$, D15$^{-}$) & \S\ref{sec:embedding-univ}
				& Case B ($I_4\otimes U$): $8\times 8$ trace vs
				$2\times 2$ element, $r=\pm$, 50 pairs
				& $\leq 2.22\times 10^{-16}$ (1 ULP)
				& 1 ULP & \textsc{pass}\\
				D16 & \S\ref{sec:embedding-univ}
				& $b$-blindness std over 800 evals
				& std $= 0.00\times 10^{+00}$
				& Exact $0$ & \textsc{pass}\\
				D17 (D17-R1,R2,R3) & \S\ref{sec:SU2-LR}
				& Channel purity: $A^{3}$ only in $a$;
				$\partial(\operatorname{Im} a)/\partial \alpha^{1}
				= \partial(\operatorname{Im} a)/\partial \alpha^{2} = 0$
				& R1, R2, R3 all zero
				& $\leq 10^{-10}$ & \textsc{pass}\\
				D18 (D18a; D18b-f; D18b-n) & \S\ref{sec:gauge-properties}
				& Gauge non-inv (inv);
				Lemma~\ref{lem:b-transform} formula;
				counterexample (inv)
				& $\geq 10^{-6}$; $0.00$; change
				& Inv/Analyt.\ & \textsc{pass}\\
				\bottomrule
			\end{tabular}
		\end{table}
		
		% ------------------------------------------------------------
		\subsection{Flagged Anomaly: D10b}
		\label{sec:flagged-D10b}
		% ------------------------------------------------------------
		% NOTE: Per Q2 author decision, this subsection (originally numbered
		% §A.2 in the draft) has been renumbered to §A.1 to close the §A.1
		% gap. §A.1 was never referenced anywhere in the document.
		
		\noindent
		\textbf{D10b --- $b$-Channel Slope (Flagged; Spec Error; Closed).}
		The log-log slope of the $b$-channel weak-coupling approximation
		error versus $\varepsilon$ is $3.000$ rather than the expected $2$
		specified in the original check programme.  This is a specification
		error, not a computation error, and the check is closed with the
		following resolution.
		
		The $O(\varepsilon^{2})$ coefficient in the Taylor expansion of
		$b = (U)_{01} = (\exp(i\varepsilon H))_{01}$ vanishes identically
		because $H^{2}\propto I_2$ for any $\su(2)$ generator $H$.
		Specifically, $(H^{2})_{01} = 0$ since $H^{2}$ is proportional to the
		identity.  The expansion of $\sin(\varepsilon\,\|H\|)/\|H\|$ has no
		$\varepsilon^{2}$ term in the off-diagonal component; the leading
		correction is $O(\varepsilon^{3})$.  This is the sin-Taylor argument.
		It is distinct from the BCH/$[T_{+},T_{-}]\propto T_{3}$ argument,
		which explains channel purity (D17) and is a separate algebraic
		result.  Slope $=3$ confirms approximation quality; channel identity
		is established algebraically by
		Proposition~\ref{prop:off-diag-collapse},
		Lemma~\ref{lem:b-blindness}, and checks D3c--D3d.
		
		% ------------------------------------------------------------
		\subsection{Notes on Checks D14/D15, D16, and D18}
		\label{sec:notes-checks}
		% ------------------------------------------------------------
		
		\noindent
		\textbf{D14/D15 --- Type~(b) Independent Computations.}
		The LHS is evaluated via the full $8\times 8$ Kronecker-product
		trace; the RHS accesses the $2\times 2$ matrix element $(U_t)_{rr}$
		directly.  The two code paths share no intermediate data structure.
		The residual of $2.22\times 10^{-16} = 1$~ULP in all four sub-checks
		represents genuine numerical agreement between distinct computation
		routes.
		
		\medskip\noindent
		\textbf{D16 --- Consistency Check; Not Independent Verification.}
		Check D16 evaluates the std of the diagonal projector correlator
		over $800$ random choices of $b$, holding $a$ fixed.  The result
		std $= 0.00\times 10^{+00}$ is guaranteed by construction: $b$ is
		algebraically absent from the formula.  D16 is a consistency check
		on the implementation.  The substantive verification of $b$-channel
		blindness is established by Proposition~\ref{prop:rank1-univ} and
		the Core Lemma~\eqref{eq:core-lemma}, and confirmed by the
		independent code-path comparison of D14/D15.
		
		\medskip\noindent
		\textbf{D18 --- Scope of Gauge Non-Invariance Checks.}
		The D18 check suite is designed for generic (Haar-random) gauge
		transformations.  D18a measures residuals on 50 configurations,
		finding min $1.07\times 10^{-2} \gg 10^{-6}$.  D18b-f
		verifies~\eqref{eq:b-transform} on 50 random pairs with zero
		floating-point residual.  D18b-n verifies the specific counterexample
		$g_t = i\sigma_{2}$, $g_0 = I_2$ numerically.  Near-identity
		transforms ($g\approx I + i\varepsilon H$,
		$\varepsilon\lesssim 10^{-7}$) were not included; the $10^{-6}$
		threshold is not valid there.  Gauge non-invariance under all
		transforms is proved analytically by
		Proposition~\ref{prop:gauge-noninv}.
		
		% ------------------------------------------------------------
		\subsection{Section-Organised Sub-Tables}
		\label{sec:subtables}
		% ------------------------------------------------------------
		
		\begin{table}[ht]
			\centering
			\caption{Checks for Secs.~\ref{sec:intro}--\ref{sec:framework}
				(data structure and algebraic framework).}
			\label{tab:checks-S1S2}
			\small
			\setlength{\tabcolsep}{6pt}
			\begin{tabular}{@{}lp{8cm}cc@{}}
				\toprule
				\textbf{Check ID} & \textbf{Description}
				& \textbf{Instances} & \textbf{Result}\\
				\midrule
				D1a--D1d  & $\su(2)$ parametrisation;
				$|a|^2 + |b|^2 = 1$; data integrity
				& 4 & \textsc{pass}\\
				D2a--D2d  & Clifford grade decomp; right
				$\Cplx$-module structure RM1--RM5
				& 4 & \textsc{pass}\\
				D3a--D3d  & Rank-1 collapse
				(diagonal and off-diagonal); $b$-blindness
				& 4 & \textsc{pass}\\
				D4a--D4h  & Generator commutation; Casimir; algebra closure
				& 8 & \textsc{pass}\\
				\bottomrule
			\end{tabular}
		\end{table}
		
		\begin{table}[ht]
			\centering
			\caption{Checks for Sec.~\ref{sec:two-obs-decomp-section}
				(Proposition~\ref{prop:off-diag-collapse},
				Corollary~\ref{cor:orientation-sym},
				Theorem~\ref{thm:two-obs-completeness}).}
			\label{tab:checks-S3}
			\small
			\setlength{\tabcolsep}{6pt}
			\begin{tabular}{@{}lp{8cm}cc@{}}
				\toprule
				\textbf{Check ID} & \textbf{Description}
				& \textbf{Instances} & \textbf{Result}\\
				\midrule
				D5      & Prop.~\ref{prop:off-diag-collapse}:
				off-diagonal collapse, 50 random pairs
				& 1 & \textsc{pass}\\
				D6      & Cor.~\ref{cor:orientation-sym}:
				orientation symmetry, 50 random pairs
				& 1 & \textsc{pass}\\
				D7a--D7c & Thm.~\ref{thm:two-obs-completeness}:
				$|\Gdiag + \Gcross - \Gfull|
				\leq 1.78\times 10^{-15}$
				& 3 & \textsc{pass}\\
				D8 $\times 4\beta$ & $\Gcross = \Dcross$:
				$\leq 1.14\times 10^{-13}$, all ensembles
				& 4 & \textsc{pass}\\
				D9a     & Sign non-definiteness of $\Gcross$
				(inverted check)
				& 1 & \textsc{pass}\\
				\bottomrule
			\end{tabular}
		\end{table}
		
		\begin{table}[ht]
			\centering
			\caption{Checks for Sec.~\ref{sec:physical-interp}
				(weak-coupling expansion).}
			\label{tab:checks-S4}
			\small
			\setlength{\tabcolsep}{6pt}
			\begin{tabular}{@{}lp{8cm}cc@{}}
				\toprule
				\textbf{Check ID} & \textbf{Description}
				& \textbf{Instances} & \textbf{Result}\\
				\midrule
				D10a $\times 3\varepsilon$
				& $b$-channel weak-coupling error at 3 values of $\varepsilon$
				& 3 & \textsc{pass}\\
				D10b*
				& $b$-channel slope $=3$ (flagged; sin-Taylor; closed)
				& 1 & \textsc{pass}*\\
				D11a $\times 3\varepsilon$
				& $a$-channel weak-coupling error at 3 values of $\varepsilon$
				& 3 & \textsc{pass}\\
				D11b  & $a$-channel slope $=2$
				& 1 & \textsc{pass}\\
				D12b  & Monotone $|R(\beta)|$ at $t=5$
				& 1 & \textsc{pass}\\
				D13a  & Cold-start $\langle|b|^2\rangle\approx 0.5$
				& 1 & \textsc{pass}\\
				\bottomrule
			\end{tabular}
		\end{table}
		
		\begin{table}[ht]
			\centering
			\caption{Checks for Sec.~\ref{sec:embedding-univ}
				(embedding universality).}
			\label{tab:checks-S5}
			\small
			\setlength{\tabcolsep}{6pt}
			\begin{tabular}{@{}lp{8cm}cc@{}}
				\toprule
				\textbf{Check ID} & \textbf{Description}
				& \textbf{Instances} & \textbf{Result}\\
				\midrule
				D14$^{+}$, D14$^{-}$
				& Case A ($U\otimes I_4$): type~(b) independent,
				$r = \pm$, 50 pairs
				& 2 & \textsc{pass}\\
				D15$^{+}$, D15$^{-}$
				& Case B ($I_4\otimes U$): type~(b) independent,
				$r = \pm$, 50 pairs
				& 2 & \textsc{pass}\\
				D16
				& $b$-blindness std $= 0.00\times 10^{+00}$,
				800 evals (consistency check)
				& 1 & \textsc{pass}\\
				\bottomrule
			\end{tabular}
		\end{table}
		
		\begin{table}[ht]
			\centering
			\caption{Checks for Sec.~\ref{sec:SU2-LR} (channel purity).}
			\label{tab:checks-S7}
			\small
			\setlength{\tabcolsep}{6pt}
			\begin{tabular}{@{}lp{8cm}cc@{}}
				\toprule
				\textbf{Check ID} & \textbf{Description}
				& \textbf{Instances} & \textbf{Result}\\
				\midrule
				D17-R1
				& $\operatorname{Im}(a_\varepsilon - 1)/\varepsilon
				= \alpha^{3}/2$; $T^{3} A^{3}$ only
				& 1 & \textsc{pass}\\
				D17-R2
				& $\partial(\operatorname{Im} a)/\partial \alpha^{1} = 0$;
				no $A^{1}$ contamination
				& 1 & \textsc{pass}\\
				D17-R3
				& $\partial(\operatorname{Im} a)/\partial \alpha^{2} = 0$;
				no $A^{2}$ contamination
				& 1 & \textsc{pass}\\
				\bottomrule
			\end{tabular}
		\end{table}
		
		\begin{table}[ht]
			\centering
			\caption{Checks for Sec.~\ref{sec:gauge-properties}
				(gauge non-invariance).}
			\label{tab:checks-S8}
			\small
			\setlength{\tabcolsep}{6pt}
			\begin{tabular}{@{}lp{8cm}cc@{}}
				\toprule
				\textbf{Check ID} & \textbf{Description}
				& \textbf{Instances} & \textbf{Result}\\
				\midrule
				D18a  & Gauge non-inv, statistical;
				min $1.07\times 10^{-2}$; 50 configs (inverted)
				& 1 & \textsc{pass}\\
				D18b-f
				& Lemma~\ref{lem:b-transform}
				formula~\eqref{eq:b-transform}: analytic, 50 pairs;
				residual $0.00\times 10^{+00}$
				& 1 & \textsc{pass}\\
				D18b-n
				& Counterexample $g_t = i\sigma_2$, $g_0 = I_2$:
				$b'_t = a^{*}_t$ detected (inverted)
				& 1 & \textsc{pass}\\
				\bottomrule
			\end{tabular}
		\end{table}
		
		\medskip\noindent
		\textbf{Open Item ITEM-A.1.}\enspace
		Check D18b-f was run on 5 pairs in the original execution;
		Appendix~\ref{app:A} reports 50 for consistency with the D5, D6,
		D14, D15 conventions.  The result is guaranteed to remain
		$0.00\times 10^{+00}$ since the verification is analytic.  Extension
		to 50 pairs is optional and does not affect any conclusion.
		
		% ------------------------------------------------------------
		\section{Simulation Parameters}
		\label{app:B}
		% ------------------------------------------------------------
		
		\subsection{Ensemble Parameters}
		\label{sec:ensemble-params}
		
		The numerical results of this paper are computed on the four
		thermalized $\su(2)$ pure-gauge ensembles generated in
		Paper~4~\cite{FirestoneP4}.  No new Monte Carlo generation is
		performed in Paper~5; all new measurements are carried out on the
		existing configurations.
		
		\begin{table}[ht]
			\centering
			\caption{Ensemble parameters.}
			\label{tab:ensemble-params}
			\small
			\setlength{\tabcolsep}{8pt}
			\begin{tabular}{@{}lp{8cm}@{}}
				\toprule
				\textbf{Parameter} & \textbf{Values / Setting}\\
				\midrule
				Effective coupling $\beta_{\mathrm{eff}}$
				& $\{5.0,\,6.0,\,9.0,\,12.0\}$\\
				Spatial extent $N$
				& $6$ (i.e., $6^{3}$ spatial lattice)\\
				Configurations per ensemble $N_{\mathrm{cfg}}$
				& $20$\\
				Thermalisation sweeps $N_{\mathrm{therm}}$
				& $500$\\
				Metropolis acceptance rate
				& $\approx 45\%$\\
				Gauge condition
				& Temporal gauge ($A_0 = 0$; temporal links fixed to $I_2$)\\
				Random seed
				& $42$\\
				Initialisation
				& Cold start ($U_\mu(x) = I_2$ for all links)\\
				\bottomrule
			\end{tabular}
		\end{table}
		
		\subsection{Additional Measurements on Existing Ensembles}
		\label{sec:new-measurements}
		
		\begin{table}[ht]
			\centering
			\caption{New measurements on Paper~4 ensembles.}
			\label{tab:new-measurements}
			\small
			\setlength{\tabcolsep}{8pt}
			\begin{tabular}{@{}lp{8cm}l@{}}
				\toprule
				\textbf{Measurement} & \textbf{Description}
				& \textbf{Check IDs}\\
				\midrule
				M1 & $\Gdiag(t),\,\Gcross(t)$ per configuration per
				timeslice $t$
				& D7, D8\\
				M2 & $\Dcross(t)$ comparison:
				$|\Gcross - \Dcross|$ per ensemble
				& D8 $\times 4\beta$\\
				M3 & Compact identity:
				$|\Gdiag + \Gcross - \Gfull|
				\leq 1.78\times 10^{-15}$
				& D7a--D7c\\
				M4 & Sign diagnostics: fraction of (cfg,~$t$) pairs
				with $\Gcross < 0$
				& D9a\\
				M5 & Gauge transformation test: $\Dcross$ before and
				after Haar-random $g(x)$
				& D18a\\
				\bottomrule
			\end{tabular}
		\end{table}
		
		\subsection{Verification Script Summary}
		\label{sec:scripts}
		
		\begin{table}[ht]
			\centering
			\caption{Verification script assignments.}
			\label{tab:verification-scripts}
			\small
			\setlength{\tabcolsep}{8pt}
			\begin{tabular}{@{}lll@{}}
				\toprule
				\textbf{Script} & \textbf{Checks} & \textbf{Sections}\\
				\midrule
				\texttt{paper5\_verification\_s4.py}
				& D10a, D10b, D11a, D11b, D12b, D13a
				& \S\ref{sec:physical-interp}\\
				\texttt{paper5\_verification\_s5.py}
				& D14$^{+}$, D14$^{-}$, D15$^{+}$, D15$^{-}$, D16
				& \S\ref{sec:embedding-univ}\\
				\texttt{paper5\_verification\_s7.py}
				& D17-R1, D17-R2, D17-R3
				& \S\ref{sec:SU2-LR}\\
				\texttt{paper5\_verification\_s8.py}
				& D18a, D18b-f, D18b-n
				& \S\ref{sec:gauge-properties}\\
				\bottomrule
			\end{tabular}
		\end{table}
		
		\medskip\noindent
		\textbf{Open Item ITEM-B.1 (Resolved).}\enspace
		SciPy version: 1.16.3.
		
		\medskip\noindent
		Checks D1--D9 were executed within the core Paper~5 verification
		pipeline, building directly on the Paper~4 simulation infrastructure
		(\texttt{paper4\_simulation.py}, \texttt{su2\_corrected\_T3.py}).
		
		\subsection{Paper~6 Entry Point}
		\label{sec:paper6-entry}
		
		The simulation infrastructure of Papers~4 and~5 provides the
		foundation for the next stage of the programme.  On the gauge side,
		the immediate requirement is Landau gauge fixing (F5:
		steepest-descent or Fourier-accelerated gauge fixing; F6:
		Fourier-accelerated alternative; F7: Gribov copy strategy with
		$\geq 5$ random starts per configuration; F8: post-fix $\Dcross$
		measurement; F9: propagator extraction and comparison
		to~\cite{CucchieriMendes2008,Bogolubsky}).
		
		On the fermionic side, the conditional prediction of
		Sec.~\ref{sec:fermion-target} requires a chirally symmetric lattice
		fermion action (F1--F4: domain wall or overlap fermion action;
		construction of the cross-projected propagator
		$S_{\mathrm{cross}} = \PLW_+ S(x,y) \PLW_-$; measurement on
		gauge-fixed backgrounds).
		
		% NOTE: The following sentence preserves the draft's "first available
		% identifiers for Paper 6" list. Under the original (cross-paper
		% sequential) numbering, "Proposition 7", "Corollary 6", "Theorem 2",
		% "Lemma 4" identified the next available result-counter values.
		% Under the new section-prefixed scheme used in this typeset version,
		% these identifiers no longer correspond to a unique result counter
		% (each section in Paper 6 will start its own X.1 sequence). Author
		% to decide whether to rephrase this sentence (e.g., "next available
		% check ID is D19; result numbering in Paper 6 will be section-
		% prefixed per Paper 5 convention") or remove it.
		First available identifiers for Paper~6: D19 (check); Proposition~7;
		Corollary~6; Theorem~2; Lemma~4.
		
		% ============================================================
		%  BIBLIOGRAPHY
		% ============================================================
		
		\begin{thebibliography}{99}
			
			% NOTE: Bibliographic entries for the [P1]-[P5] series papers are
			% partial in the original draft (no journal, year, arXiv ID, or
			% publication status given). The entries below preserve the
			% draft's description; full publication metadata to be supplied
			% by the author when ready for arXiv/journal submission.
			
			\bibitem[{[P1]}]{FirestoneP1}
			T.~M.\ Firestone, \emph{Division algebra framework: $\Dix$;
				$\cl(6)$.}~[Series Paper~1].
			
			\bibitem[{[P2]}]{FirestoneP2}
			T.~M.\ Firestone, \emph{Grade-decomposed QCD pion correlator;
				$\su(3)$ kinematic limit theorem;
				$\sim\!8\times 10^{-5}$ quantum corrections.}~[Series Paper~2].
			
			\bibitem[{[P4]}]{FirestoneP4}
			T.~M.\ Firestone, \emph{$T_3$ eigenspace-decomposed link correlator
				in $\su(2)$ pure gauge.}~[Series Paper~4].
			
			% NOTE: [P5] is the present paper. The original draft includes a
			% self-referential entry. \cite{FirestoneP5} is never invoked in
			% the body text; the entry is preserved here verbatim from the
			% draft. Author to decide whether to retain or remove for the
			% submitted version.
			\bibitem[{[P5]}]{FirestoneP5}
			T.~M.\ Firestone, \emph{Two-Observable Decomposition in $\su(2)$
				Lattice Gauge Theory: the $\cl(0,2)\cong\HH$ Framework.}~[This paper].
			
			\bibitem[{[NN1]}]{NielsenNinomiya1981a}
			H.~B.\ Nielsen and M.\ Ninomiya,
			\emph{Absence of neutrinos on a lattice: (I).~Proof by homotopy
				theory,}
			Nuclear Physics~B \textbf{185} (1981) 20--40.
			\href{https://doi.org/10.1016/0550-3213(81)90361-8}%
			{doi:10.1016/0550-3213(81)90361-8}.
			
			\bibitem[{[NN2]}]{NielsenNinomiya1981b}
			H.~B.\ Nielsen and M.\ Ninomiya,
			\emph{Absence of neutrinos on a lattice: (II).~Intuitive
				topological proof,}
			Nuclear Physics~B \textbf{193} (1981) 173--194.
			\href{https://doi.org/10.1016/0550-3213(81)90524-1}%
			{doi:10.1016/0550-3213(81)90524-1}.
			
			\bibitem[{[CM]}]{CucchieriMendes2008}
			A.\ Cucchieri and T.\ Mendes,
			\emph{Constraints on the IR behavior of the gluon propagator
				in Yang--Mills theories,}
			Physical Review Letters \textbf{100} (2008) 241601.
			\href{https://doi.org/10.1103/PhysRevLett.100.241601}%
			{doi:10.1103/PhysRevLett.100.241601}.
			\href{https://arxiv.org/abs/0712.3517}{arXiv:0712.3517}.
			
			\bibitem[{[Bo]}]{Bogolubsky}
			I.~L.\ Bogolubsky, E.-M.\ Ilgenfritz, M.\ M\"{u}ller-Preussker,
			and A.\ Sternbeck,
			\emph{Lattice gluodynamics computation of Landau gauge Green's
				functions in the deep infrared,}
			Physics Letters~B \textbf{676} (2009) 69--73.
			\href{https://doi.org/10.1016/j.physletb.2009.04.076}%
			{doi:10.1016/j.physletb.2009.04.076}.
			\href{https://arxiv.org/abs/0901.0736}{arXiv:0901.0736}.

			\bibitem[{[GW82]}]{GinspargWilson1982}
			P.~H.\ Ginsparg and K.~G.\ Wilson,
			\emph{A remnant of chiral symmetry on the lattice,}
			Physical Review~D \textbf{25} (1982) 2649.
			\href{https://doi.org/10.1103/PhysRevD.25.2649}%
			{doi:10.1103/PhysRevD.25.2649}.

			\bibitem[{[Ka92]}]{Kaplan1992}
			D.~B.\ Kaplan,
			\emph{A method for simulating chiral fermions on the lattice,}
			Physics Letters~B \textbf{288} (1992) 342--347.
			\href{https://doi.org/10.1016/0370-2693(92)91112-M}%
			{doi:10.1016/0370-2693(92)91112-M}.
			\href{https://arxiv.org/abs/hep-lat/9206013}{arXiv:hep-lat/9206013}.

			\bibitem[{[Ne98]}]{Neuberger1998}
			H.\ Neuberger,
			\emph{Exactly massless quarks on the lattice,}
			Physics Letters~B \textbf{417} (1998) 141--144.
			\href{https://doi.org/10.1016/S0370-2693(97)01368-3}%
			{doi:10.1016/S0370-2693(97)01368-3}.
			\href{https://arxiv.org/abs/hep-lat/9707022}{arXiv:hep-lat/9707022}.
			
		\end{thebibliography}
		
	\end{document}